\chapter{Technical Glossary}
\label{ap:glossary}

Table~\ref{tab:glossary} defines all acronyms and specialized terms used throughout this report.

\begingroup\small
\begin{longtable}{@{}L{3.6cm}L{10cm}@{}}
\caption{Technical glossary} \label{tab:glossary} \\
\toprule
Term / Acronym & Full Form / Definition \\
\midrule
\endfirsthead
\toprule
Term / Acronym & Full Form / Definition \\
\midrule
\endhead
AGV & Autonomous Ground Vehicle (self-propelled wheeled robot operating without continuous human teleoperation) \\
ArUco & Augmented Reality University of C\'ordoba (fiducial marker system for pose estimation) \\
\code{CALIB\_CB\_FAST\_CHECK} & OpenCV flag enabling coarse histogram pre-check in \code{findChessboardCorners}; fails on non-standard boards \\
CB & Checkerboard (high-contrast printed calibration and localization target) \\
Chi-squared ($\chi^2$) gate & Statistical outlier-rejection test comparing a measurement's Mahalanobis distance against its expected covariance; used by the estimators to reject wheel/visual measurements inconsistent with the current state \\
Curve25519 & Elliptic-curve Diffie--Hellman key-agreement function; the asymmetric primitive underlying WireGuard's key exchange \\
EKF & Extended Kalman Filter (recursive Bayesian estimator linearizing a nonlinear model at each step; MSCKF is a variant restricted to a sliding window of past poses) \\
EMA & Exponential Moving Average (first-order IIR low-pass filter for signal smoothing) \\
FIT & Florida Institute of Technology, Melbourne FL \\
FSM & Finite State Machine (formal behavioral specification with enumerable states and guarded transitions) \\
GIL & Global Interpreter Lock (CPython mutex preventing true multicore parallelism in threaded Python) \\
HMI & Human-Machine Interface (system enabling operator monitoring and control of a robotic system) \\
IIR & Infinite Impulse Response (digital filter whose current output depends on past outputs, i.e.\ recursive) \\
IMU & Inertial Measurement Unit (sensor measuring linear acceleration and angular rate) \\
LERP & Linear Interpolation (smooth blending between two values: $x(t)=x_0+t(x_1-x_0)$, $t\in[0,1]$) \\
MINS & Multisensor-aided INertial navigation System (open-source tightly-coupled EKF estimator, RPNG/University of Delaware, fusing IMU with camera and other aiding sensors in one filter, \S\ref{sec:fusionarch}) \\
MJPEG & Motion JPEG (video format where each frame is individually JPEG-compressed; no inter-frame prediction) \\
MSCKF & Multi-State Constraint Kalman Filter (EKF variant that augments the state with a sliding window of past camera poses instead of 3D landmark positions, avoiding per-feature state growth, \S\ref{sec:fusionarch}) \\
NAT & Network Address Translation (router/firewall function rewriting private addresses to a shared public one; its UDP mapping timeout motivates WireGuard's keepalive, \S\ref{sec:network}) \\
OpenVINS & Open-source MSCKF-based Visual-Inertial Navigation System (RPNG/University of Delaware; the ``VINS'' side of this report's VINS/MINS switch, ROS node \code{ov\_msckf}) \\
PnP & Perspective-n-Point (algorithm computing camera pose from $n$ 3D-2D point correspondences) \\
RAF & \code{requestAnimationFrame} (Web API for scheduling visual updates synchronized to display refresh rate) \\
ROI & Region of Interest (rectangular subset of an image processed independently of the full frame) \\
ROS & Robot Operating System (middleware framework for robotic software: pub/sub, drivers, tools) \\
\code{rosmon} & ROS process monitor (drop-in replacement for plain \code{roslaunch} node supervision, used on the robot; exposes a \code{start\_stop} service for external start/stop control, used by the watchdog's camera-recovery ladder) \\
RTB & Return To Base (autonomous navigation back to the mission start position, i.e.\ world origin) \\
RTK-GPS & Real-Time Kinematic GPS (differential correction achieving cm-level outdoor positioning) \\
SA & Situation Awareness (Endsley's three-level cognitive model: perception, comprehension, projection) \\
SB & Sub-Pixel-Befitting (OpenCV's enhanced checkerboard detector using locally adaptive binarization) \\
SE(3) & Special Euclidean group in three dimensions: the set of rigid-body transformations (rotation + translation); the pose selector's source-switch correction (Eq.~\ref{eq:se3}) is an SE(3) element \\
SLAM & Simultaneous Localization and Mapping (algorithm building a map while localizing within it) \\
TCPROS & ROS's default transport protocol: topic data flows over a plain TCP socket whose port is negotiated per connection through an XML-RPC handshake (\S\ref{sec:network}) \\
TRS & Taux de Rendement Synth\'etique, i.e.\ Overall Equipment Effectiveness (KPI used in Nicolas's prior GVS work) \\
UBSAN & Undefined Behavior Sanitizer (compiler instrumentation detecting undefined-behavior violations at runtime; surfaced the PC Wi-Fi driver fault of \S\ref{sec:network} via \code{dmesg}) \\
UGV & Unmanned Ground Vehicle (remotely operated or autonomous wheeled/tracked platform) \\
UWB & Ultra-Wideband: radio technology (3.1--10.6~GHz) enabling cm-level indoor ranging \\
VPN & Virtual Private Network (encrypted point-to-point or site-to-site network overlay; WireGuard, \S\ref{sec:network}, is the VPN implementation used in this report) \\
VSM & Value Stream Mapping (lean manufacturing process analysis tool, Nicolas's GVS background) \\
WebSocket & RFC~6455 (full-duplex TCP-based protocol for real-time bidirectional browser communication) \\
WireGuard & Kernel-space VPN protocol using Curve25519 key exchange and ChaCha20-Poly1305 encryption over a single UDP port; replaces \code{sshuttle} as this report's off-site transport (\S\ref{sec:network}) \\
XML-RPC & Remote Procedure Call protocol encoding requests/responses as XML over HTTP; used by the ROS master for node registration and by TCPROS to negotiate each topic's data port \\
\bottomrule
\end{longtable}
\endgroup

\chapter{Hardware Specifications}
\label{ap:hardware}

\section{LEO Rover Platform}

Table~\ref{tab:leo-specs} summarizes the mechanical and electrical specifications of the LEO Rover platform.

\begin{table}[htbp]
\centering
\caption{LEO Rover mechanical and electrical specifications}
\label{tab:leo-specs}
\begin{tabular}{@{}ll@{}}
\toprule
Parameter & Value \\
\midrule
Manufacturer & Husarion (Warsaw, Poland) \\
Kinematics & 4-wheel \textbf{mecanum} (FictionLab roller set); driven
  through a differential-drive odometry approximation, \emph{not} native
  omnidirectional kinematics --- see \S\ref{sec:wheels} and
  \S\ref{ap:roadmap-kinematics} \\
Drive motors & 4$\times$ BLDC with Hall-effect encoders \\
Max speed (linear) & $\sim$0.5~m/s (firmware-limited; operational limit 0.2~m/s) \\
Max speed (angular) & $\sim$1.5~rad/s (operational limit 0.5~rad/s) \\
Chassis material & Aluminium alloy, IP44 rated \\
Battery & 14.4~V Li-ion, 7800~mAh (112~Wh) \\
Nominal run time & $\sim$4~h at slow patrol speed \\
Communication & ROS Noetic over Wi-Fi; robot AP (\code{wlan\_ext}, \code{10.0.0.1}) or, since 2026-07-20, campus Wi-Fi + WireGuard tunnel for building-wide range (Ch.~\ref{ch:fusion}, \S\ref{sec:network}) \\
Firmware & AgileX \code{leo\_firmware} on Raspberry Pi CM4 (onboard) \\
Odometry & Wheel encoder ticks $\to$ \code{/firmware/wheel\_odom} at $\sim$8~Hz \\
\bottomrule
\end{tabular}
\end{table}

\section{Intel RealSense D455 Camera}

Table~\ref{tab:d455-specs} summarizes the D455 depth camera specifications relevant to this project.

\begin{table}[htbp]
\centering
\caption{Intel RealSense D455 specifications}
\label{tab:d455-specs}
\begin{tabular}{@{}ll@{}}
\toprule
Parameter & Value \\
\midrule
Depth technology & Active IR stereo (structured light projector + 2$\times$ IR imagers) \\
IR stereo baseline & 95~mm \\
Depth range & 0.4~m--6~m (nominal); 0.1~m--8~m (extended) \\
Depth resolution & $848\times480$~px \\
Depth accuracy & $<$2\% at 4~m under controlled illumination \\
RGB resolution & $1280\times720$~px (subscribed) / $640\times480$ (compressed) \\
RGB frame rate & 30~Hz maximum; \textbf{15~Hz nominal} (lab config; measured $14.9$~Hz) \\
RGB intrinsics (measured) & $f_x=380.25$, $f_y=379.26$, $c_x=320.40$, $c_y=243.72$~px at $640\times480$ \\
IR intrinsics (measured) & $f_x=f_y=386.79$, $c_x=320.06$, $c_y=236.77$~px at $640\times480$ \\
Constant used in code & $\code{CAM\_FX}=384.65$~px (see note below) \\
Interface & USB 3.2 Gen~1 (5~Gbps) \\
SDK & Intel RealSense SDK 2.0 (\code{librealsense2}) \\
ROS driver & \code{realsense2\_camera} (ROS Noetic, \code{realsense2\_camera\_manager} node) \\
\bottomrule
\end{tabular}
\end{table}

\begin{trapbox}{Do not set the colour stream to 10\,Hz}
Earlier revisions of this table listed $10$~Hz as the nominal frame rate. That
configuration \textbf{does not exist on this sensor}. The D455 supports
$6/15/30/60/90$~Hz at $640\times480$, and the attempted change on 2026-07-08
was accepted silently by the driver, which then fell back to a depth-only
default profile and \textbf{disabled every infrared stream}. The result was a
complete vision outage that persisted until the change was reverted, with no
error message at any point. See registry item~4
(Table~\ref{tab:openitems}). \textbf{Validate any profile change against
\code{rs-enumerate-devices} before deploying it.}
\end{trapbox}

\paragraph{On the $384.65$ constant, and why it differs from both rows above.}
\code{leo\_backend.py} estimates beacon range with a single focal length,
\code{CAM\_FX} $=384.65$~px, together with an optical centre placed at the
frame centre ($320.0$, $240.0$). Neither figure is the measured calibration of
either stream. It sits $1.16\,\%$ above the measured colour $f_x$ and
$0.55\,\%$ below the measured infrared $f_x$, which is consistent with its
having been derived from the infrared sensor at a time when beacon detection
ran on the infrared stream, and then carried over unchanged when the colour
stream was adopted.

The consequence is bounded and the code says so in its own comment: the
approximation is good to roughly $1$--$2\,\%$ for navigation at $\leq 3$~m,
which is well inside the other error sources in the beacon pipeline. It is
nonetheless a constant with no single defensible provenance, and $384.65$
propagates into derived results elsewhere in this report, including the LOCK
controller's loop gain and the range estimate. Treat it as a historical
working value, not as a calibration.

\begin{trapbox}{This is a different problem from the detector intrinsics}
Do not confuse the $1$--$2\,\%$ approximation described here with the defect
of \S\ref{sec:july30intrinsics}, where \code{carolus\_astrobee\_rex} runs with
$f_x = 644.12$ and $c_x = 643.47$, calibrated at $1280\times720$ and applied
to a $640\times480$ stream. That one places the principal point outside the
image and inflates every range by $69.3\,\%$. The two are unrelated in cause
and differ by nearly two orders of magnitude in effect.
\end{trapbox}

\section{Operator Workstation}

Table~\ref{tab:workstation-specs} summarizes the operator workstation configuration used for all development and validation.

\begin{table}[htbp]
\centering
\caption{Operator workstation specifications}
\label{tab:workstation-specs}
\begin{tabular}{@{}ll@{}}
\toprule
Parameter & Value \\
\midrule
OS & Ubuntu 20.04.6 LTS (Focal Fossa), kernel 5.15 \\
ROS distribution & Noetic Ninjemys (Python 3.8) \\
CPU & Intel Core i7-10700K (8 cores, 16 threads, 3.8--5.1~GHz) \\
RAM & 32~GB DDR4-3200 \\
GPU & NVIDIA GeForce RTX~2060 (6~GB VRAM; unused for vision) \\
Network & Intel Wi-Fi 6 AX200 (802.11ax); dedicated AP for robot link \\
Python packages & \code{rospy}, \code{cv2} (4.2.0), \code{numpy} (1.17.4), \code{psutil}, \code{threading}, \code{multiprocessing} \\
Browser (cockpit) & Google Chrome 124+ / Firefox 125+ (tested); any modern WebKit/Blink \\
\bottomrule
\end{tabular}
\end{table}

\chapter{Bibliography, Figure Index, and Notation}
\label{ap:biblio}

\section{Annotated Bibliography}

The following bibliography lists all works referenced in this report, organized thematically. Citations follow the IEEE reference format; the full formatted reference list appears in the References section at the end of this document. Each entry below is accompanied by a brief annotation describing its specific relevance to the LEO Rover Mission Control project.

\subsection{Robot Operating System and Middleware}

\textbf{\cite{quigley2009ros}} Quigley et al.\ (2009) introduced the Robot Operating System (ROS), the middleware framework upon which the entire LEO Rover Mission Control system is built. The paper establishes the publish-subscribe communication model, the \code{rosmaster} name resolution service, and the package management paradigm that define ROS's architecture. The \code{/mission/telemetry}, \code{/mission/command}, and \code{/cmd\_vel} topics described in Chapter~\ref{ch:methodology} directly implement the inter-node communication protocol defined in this foundational work.

\textbf{\cite{crick2012rosbridge}} The rosbridge paper introduced the bidirectional JSON-over-WebSocket protocol that enables browser-based operator interfaces to subscribe and publish to ROS topics without requiring a ROS installation on the client. The \code{rosbridge\_websocket} server used in this project is the production implementation of the protocol described by Crick et al. The generation counter pattern (\code{currentGen}) and the \code{/mission/command} topic subscription are direct implementations of the rosbridge protocol specification (version 2.0).

\textbf{\cite{openrobotics2020noetic}} This release document specifies ROS Noetic as the final Long-Term Support (LTS) release of ROS~1, supported on Ubuntu 20.04 with Python 3.8. The platform constraints documented here: particularly the Python 3.8 compatibility requirement: governed the choice of \code{multiprocessing.get\_context('spawn')} over alternatives that may behave differently under Python 3.9+. The EOL date of May 2025 motivates the ROS~2 migration discussion in Chapter~\ref{ch:conclusion}.

\textbf{\cite{openrobotics2013rosbridgeproto}} The rosbridge protocol specification defines the JSON message format for topic advertisement, subscription, publication, and service calls over WebSocket. The \code{/mission/command} message format (\code{action}, \code{mode}, \code{lin}, \code{ang} fields) and the telemetry subscription JSON are direct implementations of the rosbridge v2.0 protocol message schema.

\subsection{Computer Vision and Image Processing}

\textbf{\cite{otsu1979threshold}} Otsu's seminal 1979 paper formalized the inter-class variance maximization criterion for automatic global threshold selection, establishing the theoretical benchmark against which the LED pipeline's fixed threshold (\code{LED\_BRIGHT\_THRESH}=220) is compared. The analysis demonstrates that Otsu's method is sub-optimal for the trimodal intensity histogram produced by LED beacons under fluorescent laboratory lighting.

\textbf{\cite{niblack1986digital}} Niblack's 1986 textbook chapter on locally adaptive thresholding introduced the formula $T(x,y)=\mu(x,y)+k\sigma(x,y)$ that forms the mathematical foundation of the adaptive binarization used inside OpenCV's \code{findChessboardCornersSB} function.

\textbf{\cite{sauvola2000adaptive}} Sauvola and Pietik\"ainen (2000) refined Niblack's formula with a normalization that makes the threshold less sensitive to absolute intensity levels and more responsive to local contrast: the specific property that enables robust checkerboard corner detection under the non-uniform fluorescent illumination in the FIT laboratory.

\textbf{\cite{viola2001rapid}} While primarily known for its face detection cascade, the Viola--Jones paper's central contribution to this project is the integral image data structure, which reduces the computation of any rectangular region sum from $O(|W|)$ to $O(1)$ after an $O(N)$ preprocessing pass.

\textbf{\cite{bradski2000opencv}} The OpenCV library, documented in Bradski's 2000 publication, provides all the computer vision primitives used in the \code{leo\_backend.py} vision subprocess: \code{cv2.threshold}, \code{cv2.morphologyEx}, \code{cv2.findContours}, \code{cv2.findChessboardCornersSB}, \code{cv2.moments}, and \code{cv2.imdecode}.

\textbf{\cite{opencv2023sb}} The OpenCV 4.x documentation for \code{findChessboardCornersSB} describes the Sub-Pixel Befitting algorithm, which uses locally adaptive binarization and direct corner-finding rather than the saddle-point detector used by \code{findChessboardCorners}.

\textbf{\cite{suzuki1985topological}} Suzuki and Abe's border-following algorithm is the basis for OpenCV's \code{findContours} function used in the LED detection pipeline, tracing external boundaries in a single $O(N)$ pass.

\textbf{\cite{matheron1975random}} Matheron's 1975 monograph established mathematical morphology as a rigorous set-theoretic framework. The formal definitions of dilation, erosion, opening, and closing follow Matheron's original notation and proofs.

\textbf{\cite{serra1982image}} Serra's 1982 textbook provided the computational framework for applying Matheron's abstract morphological theory to digital images, including the discrete approximation of the Euclidean disk as an elliptical structuring element.

\subsection{Indoor Localization and Autonomous Navigation}

\textbf{\cite{thrun2005probabilistic}} Thrun, Burgard, and Fox's foundational textbook provides the theoretical framework for understanding the limitations of wheel odometry and the motivation for landmark-based localization correction.

\textbf{\cite{murartal2017orbslam2}} ORB-SLAM2, representing the state of the art for visual SLAM, was evaluated but not selected for the LEO Rover platform: its computational profile (2--4 CPU cores) would have saturated the laboratory workstation.

\textbf{\cite{garridojurado2014automatic}} Garrido-Jurado et al.\ introduced the ArUco fiducial marker system. The checkerboard landmark used in the present project is conceptually related: both rely on the Perspective-n-Point algorithm for pose estimation.

\textbf{\cite{olson2011apriltag}} AprilTag provides improved robustness over ArUco under partial occlusion and at long distances. Chapter~\ref{ch:conclusion} identifies AprilTag as a potential upgrade path for the landmark detection system.

\textbf{\cite{zhang2000flexible}} Zhang's planar calibration technique is the basis for all single-camera calibration in OpenCV, including the factory calibration of the D455's intrinsic parameters used throughout Chapter~\ref{ch:vision-math}.

\textbf{\cite{ester1996density}} Ester et al.'s DBSCAN algorithm is the formal predecessor to the density-based LED cluster anchor selection implemented in the vision pipeline.

\subsection{Signal Processing and Estimation Theory}

\textbf{\cite{kalman1960new}} Kalman's 1960 paper establishes the theoretical framework within which the EMA filter can be understood: the steady-state Kalman gain equals the EMA smoothing factor $\alpha$, providing a formal optimality interpretation.

\textbf{\cite{wiener1949extrapolation}} Wiener's seminal work established the minimum mean-squared-error estimator for stationary processes, providing the comparison baseline used to quantify the EMA filter's $\sim$29\% suboptimality at the signal bandwidth edge.

\textbf{\cite{xu2007enhanced}} Xu and Li analyzed higher-order cascade EMA implementations achieving sharper frequency cutoffs, motivating a possible future 2nd-order EMA improvement for the LED centroid pipeline.

\textbf{\cite{tukey1977exploratory}} Tukey's foundational work on robust statistics provides the theoretical basis for the choice of the 5th percentile as the obstacle detection statistic, rather than the mean or minimum.

\subsection{Control Theory and Stability Analysis}

\textbf{\cite{ziegler1942optimum}} Ziegler and Nichols' ultimate gain and ultimate period tuning method underlies the frequency-domain stability analysis of the LOCK state P-controller.

\textbf{\cite{zames1966inputoutput}} Zames' Circle Criterion provides the frequency-domain sufficient condition applied to prove that the velocity saturation nonlinearity cannot destabilize a system stable in its linear operating region.

\textbf{\cite{alur1994theory}} Alur and Dill's timed automata formalism provides the mathematical framework for formally specifying and verifying the LEO Rover's autonomous FSM with time-based guards and invariants.

\subsection{Safety Engineering and Standards}

\textbf{\cite{iec61508}} IEC~61508 is the foundational international standard for functional safety of programmable electronic systems; the SIL classification and six-layer defense-in-depth architecture both follow this standard's recommendations.

\textbf{\cite{nasa8719}} The NASA Software Safety Standard defines the framework for classifying software hazardous functions; the heartbeat deadman switch addresses the highest-risk function identified under this framework.

\textbf{\cite{isa182}} ISA-18.2 is the process industry standard for alarm system design; the cockpit's alert system implements several ISA-18.2 best practices including hysteresis bands and alarm muting.

\textbf{\cite{sobczak1962fault}} The Fault Tree Analysis technique, developed for the Minuteman ICBM program, provides the top-down deductive methodology applied to derive the combined top-level failure probability.

\subsection{Human Factors and Interface Design}

\textbf{\cite{endsley1988design}} Endsley's three-level model of situation awareness is the foundational human factors reference directly applied to the cockpit's telemetry panels, canvas overlays, and trajectory map.

\textbf{\cite{endsley1995toward}} Endsley's expanded theory identifies the factors that degrade situation awareness, motivating the alert system's negative-attention design.

\textbf{\cite{sanders1993human}} Sanders and McCormick's textbook provides the empirical data on display luminance ratios that directly motivated the cockpit's dark background design choice.

\subsection{Network Protocols and Web Engineering}

\textbf{\cite{fette2011websocket}} RFC~6455 standardized the WebSocket protocol forming the communication backbone between the rosbridge server and the browser cockpit; the framing analysis follows this RFC's frame structure.

\textbf{\cite{bianchi2000performance}} Bianchi's analytical model for IEEE~802.11 CSMA/CA performance underlies the queuing delay computation for the dedicated operator-to-robot Wi-Fi link.

\textbf{\cite{w3c2018wcag}} WCAG~2.1 defines the contrast ratio requirements verified against the cockpit's color palette.

\textbf{\cite{w3c2023backdropfilter}} The W3C's CSS Filter Effects specification defines the \code{backdrop-filter} property enabling the glassmorphic blur effect, and its compositing cost model explains the 12-panel layout limit.

\subsection{Robotics Platforms and Hardware}

\textbf{\cite{husarion2021leo}} The Husarion LEO Rover technical documentation provides the platform specifications cited throughout this report, including kinematics, encoders, battery, and firmware topic schemas.

\textbf{\cite{intel2021d455}} The D455 datasheet provides the hardware specifications used in all sensor modeling sections: stereo baseline, depth range and accuracy, RGB resolution, and calibrated focal length.

\textbf{\cite{agilex2021limo}} The AgileX LIMO Rover documentation was consulted during the Carolus comparative platform analysis contextualizing the LEO Rover work.

\subsection{Probability, Statistics, and Reliability Theory}

\textbf{\cite{kaplan1958nonparametric}} The Kaplan--Meier survival function estimator provides the non-parametric framework appropriate for refining the safety layer failure rate estimates from operational field data.

\textbf{\cite{shannon1948mathematical}} Shannon's information theory provides the theoretical basis for modeling the heartbeat protocol as a binary channel with a computable channel capacity and error-correction buffer.

\textbf{\cite{wilson1927probable}} Wilson's score confidence interval for binomial proportions is used for computing the confidence interval of the LED and checkerboard detection rates.

\subsection{3D Graphics and Web Visualization}

\textbf{\cite{cabello2021threejs}} Three.js r128 provides the WebGL-based 3D rendering infrastructure for the satellite and space capsule decorative elements in the cockpit.

\textbf{\cite{greensock2021gsap}} GSAP's ScrollTrigger plugin enables the scroll-linked rotation of the Three.js satellite model in the logbook panel.

\section{Index of Tables}

Table~\ref{tab:table-index} lists all data tables appearing in this report, with their chapter location and a brief description of their content.

\begingroup\small
\begin{longtable}{@{}L{2.2cm}L{3.6cm}L{7.4cm}@{}}
\caption{Index of all tables in this report} \label{tab:table-index} \\
\toprule
Table & Chapter/Section & Title \\
\midrule
\endfirsthead
\toprule
Table & Chapter/Section & Title \\
\midrule
\endhead
\ref{tab:intro-localization} & Introduction: I.4.2 & Comparative survey of indoor localization paradigms \\
\ref{tab:intro-latency} & Introduction: I.5.1 & End-to-end operator latency budget \\
\ref{tab:lit-ros-comparison} & Chapter~1: 1.2.2 & ROS Noetic vs.\ ROS 2 Humble feature comparison \\
\ref{tab:thread-inventory} & Chapter~2: 2.2 & \code{leo\_backend.py} daemon thread inventory \\
\ref{tab:topic-graph} & Chapter~2: 2.3 & Complete ROS topic graph \\
\ref{tab:constants-ref} & Chapter~2: 2.7 & Complete system constants reference \\
\ref{tab:autolog-events} & Chapter~2: 2.8 & AutoLogbook event taxonomy \\
\ref{tab:fsm-transitions} & Chapter~3: 3.4.1 & Complete FSM transition table \\
\ref{tab:before-after} & Chapter~3: 3.7 & System performance before/after redesign \\
\ref{tab:timing-jitter} & Chapter~3 Ext.\: 3.K.1 & Vision pipeline stage-by-stage timing analysis \\
\ref{tab:fsm-colors} & Chapter~4: 4.5.1 & FSM state banner color encoding \\
\ref{tab:safety-layers} & Chapter~5: 5.1 & Six safety layer stack \\
\ref{tab:layer-failure-rates} & Chapter~5 Ext.\: 5.A.2 & Failure rate estimation per safety layer \\
\ref{tab:vision-perf-measurements} & Chapter~6: 6.6 & Vision pipeline performance before/after \\
\ref{tab:fsm-validation} & Chapter~6: 6.7 & FSM behavioral validation test results \\
\ref{tab:ui-perf} & Chapter~6: 6.8 & Cockpit rendering performance, 3 hardware configs \\
\ref{tab:cb-scale-sweep} & Chapter~6 Ext.\: 6.E.2 & CB detection rate/latency vs.\ scale factor \\
\ref{tab:full-mission-stats} & Chapter~6 Ext.\: 6.F.1 & Full autonomous mission statistics (45~min) \\
\ref{tab:objective-achievement} & Chapter~7: 7.1 & Objective achievement summary \\
\ref{tab:glossary} & Appendix~A & Technical glossary (27 terms) \\
\ref{tab:leo-specs} & Appendix~B: B.1 & LEO Rover mechanical/electrical specifications \\
\ref{tab:d455-specs} & Appendix~B: B.2 & Intel RealSense D455 specifications \\
\ref{tab:workstation-specs} & Appendix~B: B.3 & Operator workstation specifications \\
\bottomrule
\end{longtable}
\endgroup

\section{Index of Code Listings}

Table~\ref{tab:listing-index} catalogues the source code listings appearing in this report, with their section location, programming language, and a brief description.

\begingroup\small
\begin{longtable}{@{}L{2.6cm}L{1.4cm}L{1.8cm}L{7.4cm}@{}}
\caption{Index of source code listings} \label{tab:listing-index} \\
\toprule
Listing & Section & Language & Description \\
\midrule
\endfirsthead
\toprule
Listing & Section & Language & Description \\
\midrule
\endhead
\ref{lst:spawn-context} & 2.5.1 & Python & Spawn context declaration for the vision subprocess \\
\ref{lst:queue-semantics} & 2.5.2 & Python & Frame/result queue configuration and producer logic \\
\ref{lst:start-leo} & 2.6 & Bash & \code{start\_leo.sh}: deterministic \code{ROS\_IP} + idempotent launch \\
\ref{lst:telemetry-schema} & 2.4 & Python & Selective telemetry JSON schema \\
\ref{lst:morph-close} & 3.1.1 & Python & Morphological closing with elliptical SE \\
\ref{lst:held-flag} & 3.2.2 & Python & Held-flag hysteresis mechanism \\
\ref{lst:p-controller} & 3.4.2 & Python & LOCK-state visual servoing P-controller \\
\ref{lst:yaw-wrap} & 3.4.3 & Python & Wrap-safe yaw delta accumulation \\
\ref{lst:world-pos} & 3.5.2 & Python & World position computation \\
\ref{lst:reset-continuity} & 3.5.3 & Python & Local odometry reset with world-frame continuity \\
\ref{lst:percentile-obstacle} & 3.6.2 & Python & 5th-percentile obstacle detection statistic \\
\ref{lst:color-tokens} & 4.1.2 & CSS & Color design-token system \\
\ref{lst:glassmorphism} & 4.1.3 & CSS & Glassmorphic panel implementation \\
\ref{lst:bento-grid} & 4.2 & CSS & Bento grid layout with responsive breakpoint \\
\ref{lst:raf-loop} & 4.3.1 & JavaScript & \code{requestAnimationFrame} loop structure \\
\ref{lst:dpr-handling} & 4.3.2 & JavaScript & Device pixel ratio-aware canvas resizing \\
\ref{lst:coord-mapping} & 4.3.3 & JavaScript & Letterbox-aware coordinate mapping \\
\ref{lst:lerp-smoothing} & 4.4 & JavaScript & LERP smoothing between telemetry updates \\
\ref{lst:alert-config} & 4.6.1 & JavaScript & Alert threshold configuration \\
\ref{lst:mute-persistence} & 4.6.2 & JavaScript & Alert mute persistence \\
\ref{lst:js-ring-buffer} & 4.7.1 & JavaScript & $O(1)$ ring buffer for trajectory rendering \\
\ref{lst:outlier-rejection} & 4.7.2 & JavaScript & Trajectory outlier rejection \\
\ref{lst:world-to-canvas} & 4.7.3 & JavaScript & World-frame to canvas-pixel transform \\
\ref{lst:net-latency} & 4.8 & JavaScript & Rolling telemetry latency meter \\
\ref{lst:satellite-scene} & 4.9.1 & JavaScript & Decorative satellite 3D scene geometry \\
\ref{lst:gsap-scroll} & 4.9.2 & JavaScript & GSAP ScrollTrigger-driven rotation \\
\ref{lst:i18n-system} & 4.10 & JavaScript & Flat key-value i18n dictionary system \\
\ref{lst:safe-stop} & 5.2.1 & Python & Redundant emergency-stop publication \\
\ref{lst:stop-pending} & 5.2.2 & JavaScript & Disconnection-proof emergency stop delivery \\
\ref{lst:heartbeat-fsm} & 5.3.1 & Python & Three-state heartbeat deadman \\
\ref{lst:manual-deadman} & 5.4 & Python & Manual-mode deadman timeout \\
\ref{lst:camera-watchdog} & 5.5 & Python & Camera watchdog \\
\ref{lst:velocity-clamp} & 5.6 & Python & Velocity clamping at every publish call \\
\ref{lst:autologbook-write} & 5.7 & Python & Atomic, thread-safe logbook write \\
\ref{lst:cs1-diagnosis} & 6.2.2 & Python & Diagnostic blob log (Case Study~1) \\
\ref{lst:cs1-rootcause} & 6.2.3 & Python & Original buggy single-anchor clustering \\
\ref{lst:cs1-fix} & 6.2.4 & Python & Exhaustive anchor evaluation fix \\
\ref{lst:cs4-fix} & 6.5.3 & Bash & Deterministic \code{ROS\_IP} configuration \\
\ref{lst:cs4-fallback} & 6.5.3 & Python & Fallback IP discovery via UDP routing table \\
\bottomrule
\end{longtable}
\endgroup

\section{Mathematical Notation and Symbols}

Table~\ref{tab:notation} defines the mathematical symbols and notation used consistently throughout the report.

\begingroup\small
\begin{longtable}{@{}L{2.6cm}L{7.6cm}L{2.2cm}@{}}
\caption{Mathematical notation and symbols} \label{tab:notation} \\
\toprule
Symbol & Definition & First Used \\
\midrule
\endfirsthead
\toprule
Symbol & Definition & First Used \\
\midrule
\endhead
$\alpha$ & EMA smoothing factor ($\alpha=0.45$ LED, $\alpha=0.60$ landmark) & \S3.3.1 \\
$\alpha_{\text{optimal}}$ & Kalman-optimal EMA parameter derived from SNR & \S6.C.2 \\
$A\oplus B$ & Morphological dilation of set $A$ by structuring element $B$ & \S3.B.4 \\
$A\ominus B$ & Morphological erosion of set $A$ by structuring element $B$ & \S3.B.4 \\
$A\circ B$ & Morphological opening (erosion then dilation) & \S3.B.4 \\
$A\bullet B$ & Morphological closing (dilation then erosion) & \S3.B.4 \\
$c_x, c_y$ & Camera principal point (optical axis intersection) [px] & \S3.A.1 \\
CR & WCAG color contrast ratio between foreground and background & \S4.A.2 \\
$d_{\text{est}}$ & Estimated beacon distance from thin-lens formula [m] & \S I.4.3 \\
$\delta d$ & Total distance estimation uncertainty [m] & \S3.F.2 \\
$d_s$ & Stereo disparity [px] & \S3.I.1 \\
$\text{err\_frac}$ & Normalized bearing error in LOCK state, $\in[-1,+1]$ & \S3.4.2 \\
$f_c$ & Filter $-3$~dB cutoff frequency [Hz] & \S3.3.3 \\
$f_s$ & Sampling rate of control loop or vision pipeline [Hz] & \S3.3.1 \\
$f_x, f_y$ & Camera focal lengths in pixels & \S3.A.1 \\
$g\in SE(2)$ & Rigid body transformation element of $SE(2)$ & \S5.G.1 \\
$H(z)$ & Z-domain transfer function of EMA filter & \S3.3.1 \\
$H(e^{j\omega})$ & Frequency response (DTFT) of EMA filter & \S3.3.1 \\
$K$ & Camera intrinsic matrix ($3\times3$) & \S3.A.1 \\
$K_p$ & Proportional controller gain [rad/s per px] & \S 3.G.3 \\
$K_u$ & Ziegler--Nichols ultimate gain [rad/s per px] & \S 3.G.2 \\
$\lambda$ & Component failure rate [failures/hour] & \S5.A.1 \\
$L_1, L_2$ & CIE relative luminance of foreground/background & \S4.A.2 \\
$\mu(x,y)$ & Local mean intensity in adaptive threshold window & \S3.B.2 \\
$\mu_{\ln}, \sigma_{\ln}$ & Log-normal distribution parameters & \S6.B.1 \\
$N(b_i, r)$ & Density neighborhood of blob $b_i$ at radius $r$ & \S3.D.1 \\
$\omega$ & Angular frequency [rad/s] (control) or [rad/sample] (signal) & \S3.3.1 \\
$\omega_n$ & Natural frequency of closed-loop controller [rad/s] & \S 3.G.5 \\
$\omega_{\text{pc}}$ & Phase crossover frequency [rad/s] & \S 3.G.2 \\
$\hat p$ & Observed detection rate (sample proportion) & \S6.A.1 \\
PFH & Probability of dangerous failure per hour (IEC~61508) & \S5.C.1 \\
PM & Phase margin of P-controller [degrees] & \S 3.G.4 \\
$r$ & Wheel radius of LEO Rover [m] & \S3.H.1 \\
$R\in SO(2)$ & 2D rotation matrix (Special Orthogonal group) & \S5.G.1 \\
$\rho$ & Traffic intensity (queuing theory), $\rho=\lambda/\mu$ & \S4.D.2 \\
$s_{\text{px}}$ & Horizontal pixel span of LED cluster [px] & \S I.4.3 \\
$S(x,y)$ & Integral image (prefix sum) at pixel $(x,y)$ & \S3.B.3 \\
$SE(2)$ & Special Euclidean group in 2D (rigid body transforms) & \S5.G.1 \\
SIL & Safety Integrity Level (IEC~61508 classification) & \S5.C.1 \\
SNR & $\sigma_w/\sigma_v$: signal-to-noise standard deviation ratio (Kalman) & \S6.C.2 \\
$\sigma(x,y)$ & Local standard deviation in adaptive threshold window & \S3.B.2 \\
$\sigma_n$ & Measurement noise standard deviation [px] & \S6.C.2 \\
$\sigma_w$ & Process noise standard deviation [px/sample] & \S6.C.2 \\
$T_d$ & Total vision pipeline delay [s] (plant model) & \S 3.G.1 \\
$T(x,y)$ & Adaptive threshold value at pixel $(x,y)$ & \S3.B.2 \\
$W$ & Physical width of LED beacon assembly [m] & \S I.4.3 \\
$W_q$ & Mean queuing wait time (M/D/1 model) [s] & \S4.D.2 \\
$(w_x, w_y)$ & Robot position in world frame [m] & \S3.5.1 \\
$x[n], u[n]$ & EMA filter output and input sequences & \S3.3.1 \\
$Z(w_x, w_y)$ & Stereo depth at pixel $(w_x,w_y)$ [mm or m] & \S3.I.1 \\
\bottomrule
\end{longtable}
\endgroup

% =============================================================================
\chapter{Reproduction guide for future students}
\label{ap:repro}

This appendix is written for the next student inheriting the platform, and it
assumes \emph{no} prior knowledge of this project, of the LEO~Rover, or of
ROS. It is deliberately exhaustive: every command is given in full, with the
directory it must be run from and the output you should see. Following it from
top to bottom rebuilds the entire system from a bare machine and reproduces
every measurement in this report.

Read \S\ref{ap:repro-overview} before typing anything. Ten minutes spent on
the mental model there will save hours later, because the majority of
newcomer confusion on this platform traces back to one question: \emph{which
of the two computers is doing this?}

Paths are given relative to the project root \code{\textasciitilde/TOUT}.
The target platform is ROS~Noetic on Ubuntu~20.04, built with
\code{catkin\_tools} (\code{catkin build}, \emph{not} \code{catkin\_make}).

\section*{The short version: what will cost you a week}
\addcontentsline{toc}{section}{The short version: what will cost you a week}

Read this page now and return to it whenever something behaves strangely.
Each entry below is a failure that consumed at least a day of this project,
and each shares the same property: \textbf{the symptom points at the wrong
subsystem}. They are listed in the order you are most likely to meet them.

{\small\linespread{1.06}\selectfont
\setlength{\LTleft}{0pt}\setlength{\LTright}{0pt}
\begin{longtable}{@{}L{6.3cm}L{6.6cm}r@{}}
\toprule
What you will observe & What it actually is & Where \\
\midrule
\endfirsthead
\toprule
What you will observe & What it actually is & Where \\
\midrule
\endhead
\bottomrule
\endlastfoot
The page loads perfectly, every button is dead
  & TLS port $9443$ open but detached from the ROS graph after a master
    restart. Reload, or use port $9090$
  & \S\ref{ap:repro-webtransport} \\
\addlinespace
Neither manual nor autonomous driving responds
  & the serial link, not the logic. Both modes share
    \code{/serial\_node}, sole subscriber of \code{/cmd\_vel}
  & \S\ref{ap:repro-hitrun} \\
\addlinespace
Estimators diverge, IMU ``bugs'', trajectories jump
  & the Pi is at $86\,^\circ$C and capped to $600$\,MHz, dropping IMU samples
  & \S\ref{ap:repro-thermal} \\
\addlinespace
The pose is plausible but alternates between two values
  & a TF frame with two parents. Set \code{base\_frame} to
    \code{base\_footprint}
  & \S\ref{ap:trap-tf} \\
\addlinespace
The robot over- or under-rotates consistently
  & \code{angular\_velocity\_multiplier}, an actuation gain. Fix it
    \emph{before} touching \code{gyro\_scale}
  & \S\ref{ap:trap-diffdrive} \\
\addlinespace
Beacon positions look clean, smooth and repeatable
  & and are metrically meaningless: the intrinsics belong to a different
    resolution
  & \S\ref{ap:repro-absolute} \\
\addlinespace
The stack restarts every couple of minutes
  & a \code{PATH} defect in the cron watchdog, not an instability
  & \S\ref{ap:repro-watchdog} \\
\end{longtable}}

\noindent Two habits follow from this table and are worth adopting before you
need them. First, \textbf{probe a topic that is supposed to be noisy}: a
silent topic proves nothing, and this project lost an evening to an idle
\code{/mission/log}. Second, \textbf{prefer the narrow intervention}: cycling
one node beats restarting a stack, and almost every time the surrounding
infrastructure was blamed here, it was innocent.

\section{Step $-1$: the absolute ground zero}
\label{ap:repro-groundzero}
Everything past this point assumes \code{\textasciitilde/TOUT} already exists
on your machine and that you can already reach the robot over SSH. Neither is
true on the first day, and earlier revisions of this appendix jumped straight
to \code{cd \textasciitilde/TOUT} without ever saying how it got there. This
section is the missing zero.

\subsection{Getting the code}
\label{ap:repro-clone}
\begin{lstlisting}[language=bash]
git clone <URL_OF_THIS_REPOSITORY> ~/TOUT
cd ~/TOUT
\end{lstlisting}
\code{<URL\_OF\_THIS\_REPOSITORY>} is a placeholder, not an omission: as of
this revision the repository is private. Replace it with the SSH or HTTPS
clone URL your lab gives you, in the form
\code{git@github.com:<account>/<repo>.git}. Everything from
\S\ref{ap:repro-build} onward, including \code{requirements.txt} and
\code{tools/robot\_env.sh}, arrives with this one command; nothing in this
guide requires a second download.

\subsection{Powering on and reaching the robot for the first time}
\label{ap:repro-firstboot}
\begin{trapbox}{What follows is verified project history, not a factory
manual}
Husarion's own documentation, \url{https://docs.leorover.tech}
(\cite{husarion2021leo}), is the authority on the physical unit: power
button location, battery connector, and first-boot behaviour vary across LEO
Rover hardware revisions, and no one on this project photographed or
recorded that procedure. The credentials below are not copied from that
manual either --- they are what this project's own earliest working notes
(\code{GUIDE\_DEMARRAGE\_LEO.md}, 2026-06-23, and the network log of
2026-07-20) record as having worked, before any of the later network
changes in this appendix were made. \textbf{If your unit is a different LEO
Rover, its hotspot suffix and possibly its default password will differ.}
Confirm against Husarion's documentation before assuming these are
universal.
\end{trapbox}

\begin{enumerate}[leftmargin=1.4em]
\item \textbf{Power on.} Consult
  \url{https://docs.leorover.tech} for the switch location and charging
  state indicators on your specific unit; this is not something the present
  project verified independently. Wait for the boot sequence to settle,
  on the order of a minute.
\item \textbf{Join the robot's own WiFi hotspot} from your workstation. Its
  name has the form \code{LeoRover-XXXX}, a per-unit suffix; the unit used
  throughout this project answered to \code{LeoRover-e138}
  (confirmed live, 2026-07-20). \textbf{The WPA passphrase for this hotspot
  network is not recorded anywhere in this project's history} and is the one
  genuine gap this section cannot close: check the label on the robot
  chassis, the Husarion account the unit was purchased under, or ask
  whoever last handled the physical hardware.
\item \textbf{SSH in.} Once joined, the robot answers at a fixed address on
  its own hotspot subnet:
\begin{lstlisting}[language=bash]
ssh pi@10.0.0.1
\end{lstlisting}
  The credential this project's earliest notes record as working, before any
  rotation, was password \code{raspberry} for user \code{pi} --- the
  Raspberry Pi OS default, consistent with an unmodified factory image.
  \textbf{Change it immediately} (\code{passwd}), and proceed to
  \S\ref{ap:repro-net}'s Step~4 to replace password auth with a key before
  doing anything else over this link.
\item \textbf{If the unit needs reflashing} (corrupted SD card, or a bare
  replacement board), the base OS image and flashing instructions are
  Husarion's to provide: \url{https://docs.leorover.tech}. This project
  never reflashed a unit and has no first-hand verification of that
  procedure to offer beyond the pointer.
\end{enumerate}

\noindent Once \S\ref{ap:repro-firstboot} succeeds once, everything past it
in this appendix, WireGuard, Cloudflare, the two-workspace build, applies
identically whether you are on the hotspot, on lab WiFi, or on the tunnel:
they all resolve to \emph{some} \code{ROBOT\_HOST}, and
\code{tools/robot\_env.sh} (\S\ref{ap:repro-net}) is what stops that from
mattering ever again.

\section{Start here: what this machine actually is}
\label{ap:repro-overview}
There are \textbf{two computers}, and they are not equals.

The \textbf{robot} carries a Raspberry~Pi~4. It owns the sensors and the
motors, and, this is the part that surprises people, it also runs the
\emph{ROS master}, the name server every ROS process must register with. The
Pi is deliberately kept modest: it moves data, it does not compute
trajectories.

The \textbf{workstation} (the lab PC) runs the two estimators, the web
cockpit, and every analysis tool. It is a \emph{client} of the robot's master.
If the robot is off, nothing works. If the PC is off, the robot still drives
but nothing is estimating its position.

Three consequences to memorise now, because they explain most of
\S\ref{ap:repro-trouble}:

\begin{enumerate}[leftmargin=1.4em]
\item \textbf{Every ROS command you type on the PC talks to a master on the
  robot, over WiFi.} A command that ``hangs'' is usually a network symptom,
  not a software bug.
\item \textbf{The Pi is the bottleneck, and it runs hot.} When it overheats it
  halves its own clock, and the first casualty is sensor data
  \emph{arriving on time}. Estimators then fail in ways that look like
  estimator bugs and are not (\S\ref{ap:repro-thermal}).
\item \textbf{Nothing is ground truth.} Wheel odometry, MINS and openVINS all
  estimate. They can only be compared \emph{to each other} unless you bring a
  tape measure (\S\ref{ap:repro-calibrun}).
\end{enumerate}

\subsection{The hardware chain}
\label{ap:repro-hw}
Figure~\ref{fig:apphw} traces the physical path from a wheel encoder to the
workstation. Two links on it deserve attention, because both have caused real
failures documented in this report.

\begin{figure}[htbp]
\centering
\begin{tikzpicture}[
  font=\small, x=1cm, y=1cm,
  box/.style={draw, rounded corners=2pt, align=center, inner sep=4pt},
  hw/.style ={box, fill=blue!6,  minimum width=2.4cm, minimum height=1.0cm},
  pi/.style ={box, fill=green!8, minimum width=3.0cm},
  pc/.style ={box, fill=orange!10, minimum width=3.0cm},
  lnk/.style={-{Latex[length=2.2mm]}, thick},
  badge/.style={circle, draw=red!70!black, fill=white, inner sep=1pt,
                font=\scriptsize\bfseries, text=red!70!black}
]
  \node[hw] (mot) at (0,2.7)  {4 motors\\+ encoders};
  \node[hw] (imu) at (0,1.3)  {IMU\\(inside CORE2)};
  \node[hw] (cam) at (0,-0.9) {RealSense\\D455};

  \node[box, fill=blue!6, minimum width=2.0cm, minimum height=2.4cm]
       (core) at (3.5,2.0) {CORE2\\board\\[1pt]\scriptsize firmware};

  \node[pi, minimum height=3.6cm] (pi) at (8.2,1.2)
       {\textbf{Raspberry Pi 4}\\[3pt]\scriptsize ROS master\\
        \scriptsize serial\_node\\ \scriptsize realsense driver\\
        \scriptsize image compression};

  \node[pc, minimum height=3.6cm] (pc) at (13.2,1.2)
       {\textbf{Workstation}\\[3pt]\scriptsize MINS\\ \scriptsize openVINS\\
        \scriptsize pose\_selector\\ \scriptsize web cockpit};

  \draw[lnk] (mot.east) -- (core.north west |- mot);
  \draw[lnk] (imu.east) -- (core.west |- imu);
  \draw[lnk] (core.east) -- (pi.west |- core)
    node[midway, above=2pt, align=center, font=\scriptsize]
    {UART\\\code{/dev/ttyAMA0}};
  \node[badge] at (5.85,1.55) {1};
  \draw[lnk] (cam.east) -| (pi.south)
    node[pos=0.28, below=2pt, font=\scriptsize] {USB 3.0};
  \draw[lnk] (pi.east) -- (pc.west)
    node[midway, above=2pt, align=center, font=\scriptsize]
    {WiFi\\+ WireGuard};
  \node[badge] at (10.7,0.75) {2};
\end{tikzpicture}
\caption{The physical chain. \textbf{Weak link \textcircled{\scriptsize 1}:}
the CORE2 board talks to the Pi over a \emph{GPIO serial port}
(\code{/dev/ttyAMA0}), not USB, a detail that misleads everyone, because
unplugging USB cables cannot fix it. That link has no retransmission: if the
Pi is too busy to empty the serial buffer in time, bytes are lost outright
(\S\ref{ap:repro-uart}). \textbf{Weak link \textcircled{\scriptsize 2}:} the
WiFi hop adds latency and jitter. Measured on this install: $0\%$ packet loss
but round-trip times of $3$ to $66$\,ms. ROS tolerates this because
everything is timestamped at the source; the video preview does not, which is
why the cockpit shows a \emph{high latency} badge that is honest about the
video and irrelevant to the trajectories.}
\label{fig:apphw}
\end{figure}

\subsection{The software chain}
\label{ap:repro-sw}
ROS programs communicate by publishing named streams called \emph{topics}.
Figure~\ref{fig:appsw} shows the ones that matter, left to right: raw sensors,
then cleanup, then the two estimators, then a switch that decides which
estimator the rest of the system believes.

\begin{figure}[htbp]
\centering
\begin{tikzpicture}[
  font=\small,
  n/.style={draw, rounded corners=2pt, align=center, inner sep=3pt,
            minimum height=0.85cm, font=\scriptsize},
  raw/.style={n, fill=blue!6},
  fix/.style={n, fill=yellow!12},
  est/.style={n, fill=green!10},
  dst/.style={n, fill=orange!12},
  t/.style={-{Latex[length=1.8mm]}, semithick}
]
  \node[raw] (fimu)  at (0,3.2)  {\code{/firmware/imu}\\raw IMU, $\sim$80\,Hz};
  \node[raw] (fwh)   at (0,1.7)  {\code{/firmware/wheel\_states}\\encoders};
  \node[raw] (fcam)  at (0,0.2)  {\code{/camera/infra1|2}\\stereo IR, 15\,Hz};

  \node[fix] (san)   at (3.6,3.2) {\code{imu\_sanitizer}\\despike, rescale\\[1pt]$\rightarrow$\code{/imu/data\_clean}};
  \node[fix] (whrm)  at (3.6,1.7) {\code{wheel\_remap}\\odometry};
  \node[fix] (rep)   at (3.6,0.2) {\code{republish}\\decompress};

  \node[est] (mins)  at (7.6,2.6) {\textbf{MINS}\\cam + IMU + wheels\\[1pt]$\rightarrow$\code{/mins/imu/odom}};
  \node[est] (vins)  at (7.6,0.7) {\textbf{openVINS}\\cam + IMU only\\[1pt]$\rightarrow$\code{/ov\_msckf/odomimu}};

  \node[dst] (sel)   at (11.4,1.7) {\code{pose\_selector}\\the switch};
  \node[dst] (fused) at (11.4,0.1) {\code{/robot\_pose\_fused}\\what everyone reads};

  \draw[t] (fimu) -- (san);
  \draw[t] (fwh)  -- (whrm);
  \draw[t] (fcam) -- (rep);
  \draw[t] (san.east)  -- (mins.west);
  \draw[t] (san.east)  -- (vins.west);
  \draw[t] (whrm.east) -- (mins.west);
  \draw[t] (rep.east)  -- (mins.west);
  \draw[t] (rep.east)  -- (vins.west);
  \draw[t] (mins.east) -- (sel.west);
  \draw[t] (vins.east) -- (sel.west);
  \draw[t] (sel) -- (fused);

  \node[font=\scriptsize, align=left, draw, dashed, rounded corners,
        inner sep=4pt, fill=gray!4] at (5.7,-1.5)
       {\textbf{Why two estimators?} MINS also uses the wheels, so it starts
        while standing still and cannot drift far.\\
        openVINS uses only camera + IMU, so it is the honest test of
        \emph{vision-only} localisation --- but it must be
        \emph{shaken} to start, and nothing catches it if it slips.};
\end{tikzpicture}
\caption{The topic flow. Everything upstream of the estimators runs on the
robot; everything from \code{imu\_sanitizer} rightwards runs on the
workstation. The switch is a ROS \emph{service}, not a topic: calling
\code{/pose\_selector/set\_source} with \code{true} serves MINS,
\code{false} serves openVINS.}
\label{fig:appsw}
\end{figure}

\begin{table}[htbp]
\centering\small
\caption{Which machine runs what. \code{rosnode list} shows all of it mixed
together, which is why this split is worth knowing by heart.}
\label{tab:appwho}
\begin{tabular}{@{}lll@{}}
\toprule
Machine & Started by & Main nodes \\
\midrule
Robot (Pi) & \code{leo.service} $\rightarrow$ \code{rosmon} &
  ROS master, \code{serial\_node},\\
 & & \code{firmware\_message\_converter},\\
 & & \code{realsense2\_camera}, \code{rplidar\_node} \\
\addlinespace
Workstation & \code{roslaunch leo\_navigation} &
  \code{mins\_subscribe}, \code{ov\_msckf},\\
 & \code{navigation\_master.launch} &
  \code{pose\_selector}, \code{imu\_sanitizer},\\
 & & \code{wheel\_remap}, \code{leo\_backend},\\
 & & image \code{republish} $\times$3 \\
\addlinespace
Workstation & \code{web/start\_web.sh} &
  \code{serve.py} (:8000), \code{web\_video\_server}\\
 & & (:8080), \code{rosbridge} (:9090/:9443),\\
 & & \code{cloudflared} tunnel \\
\bottomrule
\end{tabular}
\end{table}

\subsection{Map of the repository}
\label{ap:repro-map}
Everything lives under \code{\textasciitilde/TOUT}. You will spend almost all
of your time in four of these directories; the rest is history you can ignore
until you need it.

\begin{lstlisting}[language=bash]
~/TOUT/
|-- catkin_ws/                 <-- workspace 1: openVINS + our own packages
|   `-- src/
|       |-- open_vins/         upstream openVINS (ov_msckf)
|       |-- leo_navigation/    ** OUR CODE: nodes, launch files, configs **
|       |-- leo_autonomy/      beacon/patrol behaviours, AprilTag configs
|       `-- realsense-ros/     camera driver (built here, runs on the Pi)
|-- mins_ws/                   <-- workspace 2: MINS (separate build flags)
|   `-- src/MINS/
|-- carolus_ws/                <-- workspace 3: LED beacon detection
|-- tools/                     ** scripts you will use every day **
|-- web/                       the cockpit (HTML + serve.py + start_web.sh)
|-- data/trajectories/         recorded runs (.bag) and their CSV exports
|-- logs/                      every log the system writes
|-- docs/                      protocols (e.g. floor-marking calibration)
|-- report_latex/              this report
`-- leo_backend.py             ** the cockpit's brain (single large file) **
\end{lstlisting}

\noindent The four that matter: \code{catkin\_ws/src/leo\_navigation}
(all our ROS nodes), \code{tools} (all our scripts), \code{web} plus
\code{leo\_backend.py} (the cockpit), and \code{logs} (where every answer
hides).

\begin{table}[htbp]
\centering\small
\caption{The scripts in \code{tools/} you will actually use. Every one has a
long header comment explaining \emph{why} it exists: read it before
modifying the script.}
\label{tab:apptools}
\begin{tabular}{@{}ll@{}}
\toprule
Script & What it is for \\
\midrule
\code{robot\_env.sh} & sets \code{ROS\_MASTER\_URI}/\code{ROS\_IP}; sourced by
  everything \\
\code{restart\_stack.sh} & clean restart of the PC-side stack \\
\code{record\_trajectories.sh} & records a run to a rosbag (the reliable way) \\
\code{bag\_to\_csv.py} & rosbag $\rightarrow$ one CSV per estimator \\
\code{plot\_trajectories.m} & the nine-panel MINS vs openVINS figure \\
\code{compare\_tests.m} & Test~1 vs Test~2 (two separate drives) \\
\code{calib\_terrain.py} & ground-truth calibration (terminal version) \\
\code{rosbag\_daemon.sh} & launched \emph{by \code{leo\_backend.py}} to record a
  run to disk, independently of the browser. This is the capture path to use
  for anything that matters (\S\ref{ap:repro-absolute}); the cockpit's own
  buffer is cleared by any backend restart \\
\code{launch\_carolus.sh} & the beacon detector's entry point, invoked by
  \code{tools/carolus\_detect.launch}. Start the detector through the launch
  file, not by calling the script directly, or it inherits the wrong
  environment \\
\code{leo\_watchdog.sh} & the cron watchdog, read
  \S\ref{ap:repro-watchdog} first \\
\code{update\_report.sh} & compiles and publishes this report \\
\code{check\_serial\_health.sh} & UART overrun counters
  (\S\ref{ap:repro-uart}) \\
\code{log\_pose\_static.py} & records \code{/robot\_pose\_fused} to CSV while
  the rover sits still: Protocol T1.1 (\S\ref{ap:repro-t11}) \\
\code{analyze\_pose\_variance.py} & per-axis noise and linear drift rate from
  a T1.1 CSV; separates noise from unbounded divergence \\
\bottomrule
\end{tabular}
\end{table}

\section{Upstream sources: what is ours and what is not}
\label{ap:repro-repos}
For completeness, the provenance of every component and the
workspace it lives in.

Three upstream stacks are used, each in its own catkin workspace so their
build flags and dependencies stay isolated:

\begingroup\small
\begin{longtable}{@{}L{3.2cm}L{5.6cm}L{4.6cm}@{}}
\toprule
Component & Upstream & Local workspace \\
\midrule
\endhead
openVINS (\code{ov\_msckf}) & \url{https://github.com/rpng/open_vins} &
  \code{catkin\_ws/src/open\_vins} \\
MINS & \url{https://github.com/rpng/MINS} &
  \code{mins\_ws/src/MINS} \\
Carolus beacon & \url{https://github.com/albator63/carolus_ws} &
  \code{carolus\_ws} \\
Kalibr (calibration) & \url{https://github.com/ethz-asl/kalibr} &
  \code{catkin\_ws/src/kalibr} \\
\bottomrule
\end{longtable}
\endgroup

\noindent The project-specific ROS packages (\code{leo\_navigation},
\code{leo\_autonomy}, the \code{pose\_selector}, \code{wheel\_remap} and
\code{carolus\_vicon\_bridge} nodes) live alongside openVINS in
\code{catkin\_ws/src/} and are versioned with the project, not upstream.

\section{Building the workstation from a bare machine}
\label{ap:repro-build}
\emph{Why in this order:} ROS must exist before \code{catkin} can build against
it; system libraries must exist before the estimators compile; and the two
estimator workspaces must be built \emph{separately} because they disagree
about build flags.

\subsection{Step 0: Python dependencies}
\label{ap:repro-pydeps}
Do this before anything else; several failures later in this guide are really
a missing module reported three layers away from the import.
\begin{lstlisting}[language=bash]
cd ~/TOUT
pip install -r requirements.txt
sudo apt install python3-tk        # tkinter, not pip-installable
\end{lstlisting}
\code{requirements.txt} lists what each package is for, with the versions this
project was verified against. Two entries are easy to overlook:
\code{python3-tk}, which \code{leo\_dashboard.py} and
\code{leo\_tracking\_map.py} need and which \code{pip} cannot supply, and
\code{esprima}, which is not a runtime dependency at all but the JavaScript
parser used to validate \code{web/app.js} before reloading it
(\S\ref{ap:repro-editing-web}). ROS packages such as \code{rospy},
\code{tf} and \code{cv\_bridge} come from ROS Noetic and must not be
installed with \code{pip}.

\subsection{Step 1: operating system and ROS}
\begin{enumerate}[leftmargin=1.4em]
\item Install \textbf{Ubuntu 20.04 LTS}. Not 22.04: ROS~Noetic is not packaged
  for it, and every workaround costs more time than reinstalling.
\item Install ROS~Noetic Desktop-Full, following the official instructions.
  Summarised:
\begin{lstlisting}[language=bash]
sudo sh -c 'echo "deb http://packages.ros.org/ros/ubuntu $(lsb_release -sc) main" \
  > /etc/apt/sources.list.d/ros-latest.list'
sudo apt install curl
curl -s https://raw.githubusercontent.com/ros/rosdistro/master/ros.asc | sudo apt-key add -
sudo apt update
sudo apt install ros-noetic-desktop-full
\end{lstlisting}
\item Make ROS available in every new terminal:
\begin{lstlisting}[language=bash]
echo "source /opt/ros/noetic/setup.bash" >> ~/.bashrc
source ~/.bashrc
\end{lstlisting}
  \textbf{Verify:} \code{echo \$ROS\_DISTRO} prints \code{noetic}.
\item Install the build tool and the dependency resolver:
\begin{lstlisting}[language=bash]
sudo apt install python3-catkin-tools python3-rosdep python3-osrf-pycommon
sudo rosdep init && rosdep update
\end{lstlisting}
\end{enumerate}

\subsection{Step 2: system libraries the estimators need}
Both openVINS and MINS are C++ and need Eigen, Ceres, OpenCV~4 and Boost.
\code{rosdep} finds most of them; Ceres is the one that usually needs help.
\begin{lstlisting}[language=bash]
sudo apt install libeigen3-dev libboost-all-dev libopencv-dev \
                 libceres-dev libsuitesparse-dev
# ROS-side extras used by the cockpit and the pipeline:
sudo apt install ros-noetic-rosbridge-suite ros-noetic-web-video-server \
                 ros-noetic-image-transport-plugins ros-noetic-topic-tools \
                 ros-noetic-rosmon ros-noetic-apriltag-ros
# analysis tools:
sudo apt install python3-pip texlive-full
pip3 install --user numpy pyyaml psutil pillow
\end{lstlisting}
\textbf{Note:} \code{texlive-full} is $\sim5$\,GB but avoids a long tail of
missing-package errors when compiling this report. \code{psutil} is what lets
the cockpit show CPU load; without it the cockpit still works but that panel
stays empty.

\subsection{Step 3: the three workspaces}
\emph{Why three:} openVINS and MINS come from the same research group but do
not build with the same flags, and mixing them in one workspace produces link
errors that look like source bugs. The beacon stack is separate again because
it is optional.

\begin{enumerate}[leftmargin=1.4em]
\item Create the tree and clone the upstream sources:
\begin{lstlisting}[language=bash]
mkdir -p ~/TOUT/catkin_ws/src ~/TOUT/mins_ws/src ~/TOUT/carolus_ws/src
cd ~/TOUT/catkin_ws/src
git clone https://github.com/rpng/open_vins.git
git clone https://github.com/IntelRealSense/realsense-ros.git -b ros1-legacy
cd ~/TOUT/mins_ws/src
git clone https://github.com/rpng/MINS.git
\end{lstlisting}
\item \textbf{Our packages arrive with the clone of
  \S\ref{ap:repro-clone}}, already in place at
  \code{\textasciitilde/TOUT/catkin\_ws/src/\{leo\_navigation,leo\_autonomy\}}
  --- there is nothing further to copy. They contain every node described in
  \S\ref{ap:repro-ourcode}.
  \textbf{Note on the QR scripts:} \code{qr\_leo\_control.py},
  \code{qr\_search\_move.py} and \code{diag\_qr.py} are \emph{not} a catkin
  package despite the name similarity; they are standalone Python scripts at
  the repository root, run directly with \code{python3}, not through
  \code{roslaunch} or \code{rosrun}. An earlier revision of this appendix
  referred to a \code{leo\_qr\_control} package that does not exist in the
  versioned tree; if you were looking for it, this is the correction.
\item Resolve dependencies, once per workspace:
\begin{lstlisting}[language=bash]
cd ~/TOUT/catkin_ws && rosdep install --from-paths src --ignore-src -r -y
cd ~/TOUT/mins_ws   && rosdep install --from-paths src --ignore-src -r -y
\end{lstlisting}
\item Build workspace 1 (openVINS + our packages):
\begin{lstlisting}[language=bash]
cd ~/TOUT/catkin_ws
catkin build
\end{lstlisting}
  \textbf{Expected:} a progress table ending with a green summary,
  \code{[build] Summary: All ... packages succeeded!}. On a four-core machine
  budget $15$--$25$ minutes. If \code{ov\_msckf} fails on Ceres, install
  \code{libceres-dev} and rebuild only that package with
  \code{catkin build ov\_msckf}.
\item Build workspace 2 (MINS):
\begin{lstlisting}[language=bash]
cd ~/TOUT/mins_ws
catkin build mins
\end{lstlisting}
\item Build workspace 3 (beacon, optional):
\begin{lstlisting}[language=bash]
cd ~/TOUT/carolus_ws && catkin build
\end{lstlisting}
\end{enumerate}

\subsection{Step 4: the sourcing trap, and why a wrapper exists}
\label{ap:repro-sourcetrap}
\textbf{This is the single most confusing thing about the build.} A
\code{devel/setup.bash} \emph{replaces} several environment variables rather
than appending to them. Source \code{catkin\_ws} then \code{mins\_ws}, and
\code{roslaunch} can no longer find \code{leo\_navigation}; source them the
other way round and it cannot find \code{mins}. Either way the error is
\code{cannot locate node of type [...] in package [...]}, which reads like a
missing file.

The fix is \code{catkin\_ws/src/leo\_navigation/launch\_mins.sh}: it saves the
five variables that matter, sources the second workspace, then re-appends the
saved values so both are reachable.

\begin{lstlisting}[language=bash]
source /home/lab272/TOUT/tools/robot_env.sh
source /opt/ros/noetic/setup.bash
source /home/lab272/TOUT/catkin_ws/devel/setup.bash
_RPP="$ROS_PACKAGE_PATH"; _CMPP="$CMAKE_PREFIX_PATH"; _PATH="$PATH"
_LDLP="$LD_LIBRARY_PATH"; _PYP="$PYTHONPATH"
source /home/lab272/TOUT/mins_ws/devel/setup.bash
export ROS_PACKAGE_PATH="$ROS_PACKAGE_PATH:$_RPP"
export CMAKE_PREFIX_PATH="$CMAKE_PREFIX_PATH:$_CMPP"
export PATH="$PATH:$_PATH"
export LD_LIBRARY_PATH="$LD_LIBRARY_PATH:$_LDLP"
export PYTHONPATH="$PYTHONPATH:$_PYP"
exec roslaunch leo_navigation "${1:-mins.launch}"
\end{lstlisting}

\noindent \textbf{Rule: never start anything MINS-side with a plain
\code{roslaunch}.} Always:
\begin{lstlisting}[language=bash]
~/TOUT/catkin_ws/src/leo_navigation/launch_mins.sh navigation_master.launch
\end{lstlisting}
\code{tools/restart\_stack.sh} does this for you, which is one reason to
launch through it rather than by hand.

\subsection{Quick reference: the clone-and-build commands}
\label{ap:repro-build-ref}
Clone each stack into its workspace's \code{src/}, then build. openVINS and
the project packages share \code{catkin\_ws}:

\begin{lstlisting}[language=bash]
# --- openVINS + project packages (catkin_ws) ---
cd ~/TOUT/catkin_ws/src
git clone https://github.com/rpng/open_vins.git
# (leo_navigation, leo_autonomy, kalibr already present in src/)
cd ~/TOUT/catkin_ws
catkin build ov_msckf          # builds openVINS and its deps
source devel/setup.bash

# --- MINS (separate workspace: its own build flags) ---
cd ~/TOUT/mins_ws/src
git clone https://github.com/rpng/MINS.git
cd ~/TOUT/mins_ws
catkin build mins
source devel/setup.bash

# --- Carolus beacon detection ---
cd ~/TOUT/carolus_ws
catkin build
\end{lstlisting}

\noindent System dependencies (Eigen, Ceres, OpenCV~4, Boost) are the
standard openVINS/MINS prerequisites; on a fresh machine install them with
\code{rosdep install --from-paths src --ignore-src -r -y} run once per
workspace before the first \code{catkin build}. openVINS and MINS both
compile against \code{tf2\_ros} on this platform (the legacy \code{tf}
migration of \S\ref{sec:july22tf2} for MINS, and the equivalent open item
for openVINS in Table~\ref{tab:openitems}). The MINS workspace cannot simply
be \code{source}d back-to-back with \code{catkin\_ws}; the
\code{tools/launch\_mins.sh} helper unions the two environments correctly and
is the supported way to launch anything MINS-side (see its header comment for
why).

\section{Network configuration}
\label{ap:repro-net}
\emph{Why this matters:} ROS has no auto-discovery. Each process must be told
where the master lives (\code{ROS\_MASTER\_URI}) and which address to advertise
so others can call it back (\code{ROS\_IP}). Get \code{ROS\_IP} wrong and you
get the most misleading symptom on the platform: \code{rosnode list} shows
every node, and no data ever arrives.

\subsection{Step 1: one file holds the robot's address}
The robot's address changes with the network: \code{10.0.0.1} on the robot's own
hotspot, a DHCP address on campus WiFi, \code{10.200.0.2} over the WireGuard
tunnel used at the time of writing. \textbf{Never hard-code it}: it used to
be duplicated in five scripts, and changing networks meant five edits and a
guaranteed omission.

\begin{lstlisting}[language=bash]
echo 'LEO_ROBOT_HOST=10.200.0.2' > ~/.leo_network
\end{lstlisting}
Deleting the file restores the default (\code{10.0.0.1}). Every script reads
\code{tools/robot\_env.sh}, which reads this file.

\subsection{Step 2: what robot\_env.sh derives, and why}
\code{ROS\_IP} is \textbf{not} hard-coded either: it is looked up from the
routing table, so it stays correct on a machine with both Ethernet and WiFi.
The address that matters is the one on the interface that can actually reach
the robot.
\begin{lstlisting}[language=bash]
export ROBOT_HOST="${LEO_ROBOT_HOST:-10.0.0.1}"
export ROS_MASTER_URI="http://${ROBOT_HOST}:11311"
_leo_ros_ip="$(ip route get "$ROBOT_HOST" | grep -oP 'src \K[0-9.]+' | head -1)"
export ROS_IP="${_leo_ros_ip:-$(hostname -I | awk '{print $1}')}"
\end{lstlisting}

\noindent \textbf{Verify, from any directory:}
\begin{lstlisting}[language=bash]
cd ~/TOUT
source tools/robot_env.sh
echo "$ROBOT_HOST | $ROS_MASTER_URI | $ROS_IP"
\end{lstlisting}
\textbf{Expected:} \code{10.200.0.2 | http://10.200.0.2:11311 | 10.200.0.1}. \\
If \code{ROS\_IP} is empty or on the wrong subnet, stop: fix the network
first, nothing downstream will work.

\subsection{Step 3: the WireGuard tunnel (optional but recommended)}
Campus WiFi hands out changing addresses and sometimes isolates clients from
each other. A WireGuard tunnel gives the pair two \emph{fixed} addresses
(\code{10.200.0.1} for the PC, \code{10.200.0.2} for the robot) regardless of
the underlying network, which removes an entire class of ``it worked
yesterday'' failures.

\paragraph{Decide which machine is reachable first.} WireGuard has no
client and server roles in the protocol, but it does in practice: one side
must carry an \code{Endpoint}, and that side must be reachable from the
other. This is the question that decides whether the tunnel comes up at all,
and it is decided by your network, not by preference.

\begin{center}
\small
\begin{tabular}{@{}L{4.6cm}L{4.4cm}L{4.4cm}@{}}
\toprule
Situation & Who carries \code{Endpoint} & Works? \\
\midrule
The PC has a fixed, routable address (wired lab network, port
$51820/\mathrm{udp}$ open) & the robot points at the PC & yes, the
arrangement used here \\
\addlinespace
Both sides behind campus NAT, no inbound ports & neither can reach the other
& \textbf{no}; you need a third host \\
\addlinespace
A cheap VPS with a public IP & both point at the VPS & yes, and it survives
either side changing network \\
\bottomrule
\end{tabular}
\end{center}

\noindent Campus networks routinely block inbound connections and isolate
wireless clients from one another. If the PC is itself on campus WiFi, expect
the second row: no amount of configuration on the two machines will fix it,
and the honest answer is a VPS acting as a rendezvous point, with both rover
and workstation as clients of it.

\begin{enumerate}[leftmargin=1.4em]
\item \code{sudo apt install wireguard} on both machines.
\item Generate a key pair on each:
\begin{lstlisting}[language=bash]
wg genkey | tee private.key | wg pubkey > public.key
\end{lstlisting}
\item Write \code{/etc/wireguard/wg0.conf} on each side. Substitute the four
  placeholders; everything else is literal.

\textbf{Workstation} (the side carrying the \code{Endpoint} in this
deployment, i.e.\ the one that listens):
\begin{lstlisting}[language=bash]
[Interface]
Address    = 10.200.0.1/24
PrivateKey = <PC_PRIVATE_KEY>
ListenPort = 51820

[Peer]                              # the robot
PublicKey  = <ROBOT_PUBLIC_KEY>
AllowedIPs = 10.200.0.2/32
\end{lstlisting}

\textbf{Robot:}
\begin{lstlisting}[language=bash]
[Interface]
Address    = 10.200.0.2/24
PrivateKey = <ROBOT_PRIVATE_KEY>

[Peer]                              # the workstation
PublicKey  = <PC_PUBLIC_KEY>
AllowedIPs = 10.200.0.0/24
Endpoint   = <PC_REACHABLE_IP>:51820
PersistentKeepalive = 25
\end{lstlisting}

\item Enable on both: \code{sudo systemctl enable --now wg-quick@wg0}.
\item \textbf{Verify:} \code{ping -c 5 10.200.0.2} from the PC gives
  \code{0\% packet loss}, and \code{sudo wg show} lists a recent handshake.
  Then point the project at it:
  \code{echo 'LEO\_ROBOT\_HOST=10.200.0.2' > \textasciitilde/.leo\_network}.
\end{enumerate}

\begin{trapbox}{\texttt{PersistentKeepalive} is not a roaming fix}
It is frequently described as one, including in earlier revisions of this
appendix, and the claim does not survive contact with this platform.
Keepalive stops a NAT mapping from expiring while nothing is being sent. It
does nothing when the robot's radio actually re-associates and its underlying
address changes: the peer's cached endpoint is then stale until a packet from
the new address re-teaches it.

Registry item~12 records the reality. One WireGuard outage was explained and
fixed by endpoint roaming; a \emph{second} resisted restarts on both ends, a
port change, and a four-minute idle wait, while ICMP and SSH to the robot
stayed healthy throughout. That item is still open. Treat the tunnel as
usually reliable and occasionally not, and when it floats, check
\code{sudo wg show} for the age of the last handshake before restarting
anything else.
\end{trapbox}

\noindent \textbf{Known quirk:} on this workstation the USB WiFi dongle
(\code{rtl88x2bu}) drops packets when its power-saving kicks in. It affects the
PC only, not the robot. First find your own interface name, which will
\emph{not} be the one below:
\begin{lstlisting}[language=bash]
ip -br link | grep -E '^wl'        # or: iwconfig 2>/dev/null | grep -B1 IEEE
\end{lstlisting}
Then disable power saving on it. This does \emph{not} survive a reboot:
\begin{lstlisting}[language=bash]
sudo iw dev <WIFI_INTERFACE> set power_save off
# on the machine used here that name was wlxcc641aee985b
\end{lstlisting}

\subsection{Step 4: passwordless SSH to the robot}
Several scripts read robot state over SSH. Set up a key so they do not stall
waiting for a password:
\begin{lstlisting}[language=bash]
ssh-keygen -t ed25519           # accept the defaults
ssh-copy-id pi@10.200.0.2       # password once; never again
ssh pi@10.200.0.2 'hostname'    # verify
\end{lstlisting}
\textbf{Expected:} the robot's hostname, with no password prompt.

\subsection{Step 5: the web transport, and the two ports it can use}
\label{ap:repro-webtransport}
The cockpit is an ordinary static web page. It has no server-side logic and no
ROS libraries of its own: everything it knows about the robot arrives over a
single WebSocket to \code{rosbridge}, which translates JSON into ROS messages
and back. Three processes therefore have to be alive, and they are independent
of one another:

\begin{center}
\small
\begin{tabular}{@{}llL{5.6cm}@{}}
\toprule
Process & Port & Role \\
\midrule
\code{web/serve.py}                      & $8000$ & serves the HTML, CSS and JavaScript \\
\code{rosbridge\_websocket}              & $9090$ & plain WebSocket, LAN and loopback \\
\code{tools/rosbridge\_websocket\_tunnel.py} & $9443$ & TLS WebSocket, behind \code{cloudflared} \\
\bottomrule
\end{tabular}
\end{center}

\noindent The page itself is served on $8000$ in both cases. What changes
between the two transports is only where the \emph{WebSocket} goes, and the
choice is made automatically from the hostname you typed:

\begin{lstlisting}[language=JavaScript]
const TUNNEL_ROOT = 'leo-rover-gardon.dev';
const isTunnel = () => location.hostname === TUNNEL_ROOT
                    || location.hostname.endsWith('.' + TUNNEL_ROOT);
const wsURL = (host) => isTunnel()
      ? ('wss://' + location.host + '/rosbridge')   // -> cloudflared -> :9443
      : ('ws://'  + host + ':9090');                // -> direct
\end{lstlisting}

\noindent This single function is the whole of the routing logic, and it has a
consequence worth internalising before you debug anything:

\begin{center}
\begin{tikzpicture}[font=\scriptsize, >=Latex, node distance=6mm,
  b/.style={draw, rounded corners=2pt, inner sep=3.5pt, align=center, minimum height=7mm}]
  \node[b, fill=accentblue!7, draw=accentblue] (b1) at (0,0) {browser on\\\code{localhost:8000}};
  \node[b, right=16mm of b1] (r1) {\code{rosbridge}\\\code{:9090}};
  \node[b, right=13mm of r1] (m1) {ROS master};
  \draw[->, accentblue, thick] (b1) -- node[above]{\code{ws://}} (r1);
  \draw[->] (r1) -- (m1);

  \node[b, fill=alertred!6, draw=alertred] (b2) at (0,-2.1) {browser on\\\code{cockpit.\ldots dev}};
  \node[b, right=16mm of b2] (cf) {\code{cloudflared}};
  \node[b, right=7mm of cf] (r2) {TLS bridge\\\code{:9443}};
  \node[b, right=7mm of r2] (m2) {ROS master};
  \draw[->, alertred, thick] (b2) -- node[above]{\code{wss://}} (cf);
  \draw[->] (cf) -- (r2);
  \draw[->] (r2) -- (m2);
\end{tikzpicture}
\end{center}

\noindent \textbf{The two paths share nothing but the ROS master.} Different
process, different port, different TLS state, different failure modes. When
one is broken the other is very often fine, and that is the single most useful
fact in this subsection.

\subsubsection{The zombie port}
\begin{trapbox}{Symptom: the page loads perfectly and every button is dead}
Restart the ROS master while the TLS bridge on $9443$ is running, and that
bridge keeps its listening socket open. It still accepts the TCP connection,
still completes the TLS negotiation, still completes the WebSocket upgrade.
What it no longer holds are live handles to the ROS graph, so nothing it
receives is ever published and nothing from the graph ever reaches you.

The browser has no way to notice. At every layer it can observe, the
connection succeeded. The page therefore renders its ``connected'' state
truthfully and the interface looks entirely normal while being inert.
\end{trapbox}

\noindent\textbf{Since 2026-07-31 this is detected automatically.}
\code{tools/leo\_watchdog.sh} carries a \emph{drowning probe}: it does not ask
whether the port is bound, which in this failure mode it always is, but opens
a real WebSocket handshake and requires an answer within eight seconds. Until
July~31 that probe covered only port $9090$, so the one failure mode it was
designed to catch was being observed on the port it did not watch. It now
probes both, and purges the pair on either failure, after which
\code{start\_web.sh} restarts them idempotently.

\begin{trapbox}{If you ever extend this probe, mind the TLS}
Port $9443$ is served over TLS (\code{listenSSL} in
\code{rosbridge\_websocket\_tunnel.py}). Copying the plain-text $9090$ probe
onto it produces a probe that \emph{always} fails, so the watchdog would purge
both bridges every single minute: the exact restart storm of
\S\ref{ap:repro-watchdog}. Verified on 2026-07-31, a plain probe against
$9443$ raises \code{AssertionError} while a TLS-wrapped one returns
\code{HTTP/1.1 400}. A $400$ is a perfectly good result here, because the
probe tests whether the process still answers, not whether it accepts the
request.
\end{trapbox}

\noindent\textbf{Manual remediation, in increasing order of cost}, for the
window before the watchdog's next tick:
\begin{enumerate}[leftmargin=1.6em]
\item \textbf{Force-reload the page} (\code{Ctrl+Shift+R}). This tears down
  the WebSocket and builds a new one. If the bridge process itself recovered,
  this is enough, and it costs two seconds.
\item \textbf{Change transport.} Open
  \code{http://localhost:8000/ops.html} instead of the tunnel domain. By the
  rule above this switches you to \code{ws://\ldots:9090}, a different process
  that never went through the restart. Note carefully that this is not a
  different page and not a fallback mode: it is the same cockpit reaching the
  same master by an independent route. If the cockpit works here, the fault is
  in the tunnel path and nowhere else.
\item \textbf{Restart the TLS bridge} only once the two steps above have
  failed, since it is the only one that interrupts anything.
\end{enumerate}

\noindent\textbf{Diagnose before you restart.} The two commands below separate
``the page is broken'' from ``the bus is quiet'', and the order matters:

\begin{lstlisting}[language=bash]
# 1. Does ROS traffic exist at all? Use a topic that is ALWAYS busy.
rostopic hz /mission/telemetry      # expect ~4 Hz, continuously
# 2. Does a command sent from the page reach the bus?
rostopic echo /mission/command      # press a cockpit button and watch
# 3. Are both bridges actually listening?
ss -ltnp | grep -E '9090|9443|8000'
\end{lstlisting}

\noindent The choice of topic in the first command is not incidental. During
the July 30 session the relay was declared broken on the evidence of a silent
\code{/mission/log}, which publishes only when something happens; the empty
output proved nothing at all. Repeating the test on \code{/mission/telemetry}
showed fifty frames in twelve seconds and the relay was healthy throughout.
\textbf{A silent topic is not a dead link.} Always probe something that is
supposed to be noisy.

\subsubsection{Two further transport traps, for completeness}
\begin{itemize}[leftmargin=1.4em]
\item \textbf{Opening the page as \code{file://}} disables both rosbridge and
  the video stream, because the browser refuses cross-origin WebSockets from a
  file URL. The cockpit detects this and prints a red banner telling you to
  use \code{http://<PC-IP>:8000/ops.html}. If you see that banner, nothing
  else you observe about connectivity means anything.
\end{itemize}

\subsection{Step 6: the Cloudflare tunnel, and why you must rebuild it}
\label{ap:repro-cloudflare}
Everything above gets you a cockpit on the local network. The public address,
\code{cockpit.leo-rover-gardon.dev}, is served by a Cloudflare tunnel, and it
is the one piece of this system you \textbf{cannot inherit}.

\begin{trapbox}{The existing tunnel dies with its owner's account}
The tunnel in service at the time of writing is bound to a personal
Cloudflare account and to a domain registered under it. Its credentials live
in \code{\textasciitilde/.cloudflared/} on the workstation, outside the
project tree and outside version control, and they cannot be transferred by
copying a file. \textbf{You will have to create your own tunnel, on your own
account, and repoint DNS at it.} Budget an hour. Nothing else in this
appendix depends on it: the cockpit works fully over
\code{http://localhost:8000} without any tunnel at all.
\end{trapbox}

\paragraph{What the tunnel actually does.} It terminates TLS at Cloudflare's
edge and forwards three different paths of one hostname to three local
services. That routing is the missing link between the JavaScript of
\S\ref{ap:repro-webtransport}, which builds
\code{wss://cockpit.\ldots/rosbridge}, and the TLS bridge listening on
$9443$. Without the \code{ingress} block below, that URL resolves to nothing
and every button in the cockpit is inert with no error message.

\paragraph{Setup, from nothing.}
\begin{enumerate}[leftmargin=1.4em]
\item Install the binary (a \code{.deb} from Cloudflare, or your package
  manager), then authenticate. This opens a browser and writes
  \code{\textasciitilde/.cloudflared/cert.pem}:
\begin{lstlisting}[language=bash]
cloudflared tunnel login
\end{lstlisting}
\item Create the tunnel. This prints a UUID and writes the credentials file
  \code{\textasciitilde/.cloudflared/<UUID>.json}, which is the actual
  secret:
\begin{lstlisting}[language=bash]
cloudflared tunnel create leo-cockpit
\end{lstlisting}
\item Point a DNS name at it, on a domain your account controls:
\begin{lstlisting}[language=bash]
cloudflared tunnel route dns leo-cockpit cockpit.example.org
\end{lstlisting}
\item Write \code{\textasciitilde/.cloudflared/config.yml}. This is the file
  that matters, reproduced with the real routing and only the identifiers
  replaced:
\begin{lstlisting}[language=bash]
tunnel: <TUNNEL_UUID>
credentials-file: /home/<user>/.cloudflared/<TUNNEL_UUID>.json

ingress:
  # rosbridge and video are served on the SAME hostname as the page, on
  # dedicated paths, so they share the Cloudflare Access cookie of the
  # already-authenticated page. A separate subdomain would break the
  # WebSocket handshake on cookie scope.
  - hostname: cockpit.example.org
    path: ^/rosbridge/?$
    service: https://localhost:9443
    originRequest:
      noTLSVerify: true    # self-signed cert, loopback hop only
  - hostname: cockpit.example.org
    path: ^/(stream|stream_viewer)$
    service: http://localhost:8080
  - hostname: cockpit.example.org
    service: http://localhost:8000
  - service: http_status:404
\end{lstlisting}
\item Run it. \code{web/start\_web.sh} already does this for you, with an
  absolute path because a bare \code{cloudflared} fails silently under
  \code{cron}:
\begin{lstlisting}[language=bash]
cloudflared tunnel run
\end{lstlisting}
\end{enumerate}

\paragraph{Three details that are load-bearing.}
\begin{itemize}[leftmargin=1.4em]
\item \code{service: https://localhost:9443}, not \code{http}. The TLS bridge
  speaks TLS on the loopback hop as well.
\item \code{noTLSVerify: true} is required because that bridge presents a
  self-signed certificate. It is safe \emph{here} and only here, because the
  hop it applies to never leaves the machine.
\item The order of \code{ingress} rules is significant. Cloudflare takes the
  first match, so the catch-all \code{service: http://localhost:8000} must
  come last, and \code{http\_status:404} last of all.
\end{itemize}

\paragraph{Verify.}
\begin{lstlisting}[language=bash]
pgrep -x cloudflared                            # the process must exist
curl -sI https://cockpit.example.org/ops.html   # expect HTTP/2 200
tail -20 ~/TOUT/logs/cloudflared.log            # look for "Registered tunnel"
\end{lstlisting}

\section{Robot-side configuration}
\label{ap:repro-robot}
The Pi runs a stock LEO~Rover image plus three project-specific changes.
\textbf{If you reflash the robot, you must redo all three}, and the system will
appear to work while behaving badly if you forget, which is exactly how
each was discovered.

\subsection{How the robot starts itself}
The Pi runs \code{leo.service}, a systemd unit that launches \code{rosmon}
(a supervising alternative to \code{roslaunch}), which in turn starts the ROS
master and the robot's nodes from
\code{/opt/ros/noetic/share/leo\_bringup/launch/leo\_bringup.launch}.

\begin{lstlisting}[language=bash]
ssh pi@$ROBOT_HOST 'systemctl status leo --no-pager | head -5'
ssh pi@$ROBOT_HOST 'sudo systemctl restart leo'     # full robot-side restart
\end{lstlisting}
\code{rosmon} respawns any node that dies, which is usually helpful and
occasionally the reason something you killed keeps coming back
(\S\ref{ap:repro-phantom}).

\subsection{Change 1: real-time priority for the serial node}
\label{ap:repro-rtfix}
\emph{Why:} the CORE2 board streams IMU and encoder data into a $32$-byte
hardware FIFO on the Pi's PL011 UART. If \code{serial\_node} is not scheduled
before that FIFO fills, bytes are lost with no retransmission. Symptom:
thousands of checksum errors and wheel odometry collapsing from $19$\,Hz to
$1$\,Hz.

\textbf{The fix has two halves and needs both.} With only the first, the node
does not start \emph{at all}: you lose driving entirely, which is worse than
the bug.

\begin{enumerate}[leftmargin=1.4em]
\item \textbf{Ask for real-time priority}, in
  \code{leo\_bringup.launch} on the robot:
\begin{lstlisting}[language=xml]
<node name="serial_node" pkg="rosserial_python" type="serial_node.py"
      launch-prefix="chrt -f 25">
\end{lstlisting}
\item \textbf{Grant permission to ask}, because user \code{pi} has
  \code{rtprio=0} by default and \code{chrt} would simply fail. Create
  \code{/etc/systemd/system/leo.service.d/10-rtprio.conf}:
\begin{lstlisting}
[Service]
LimitRTPRIO=50
\end{lstlisting}
  then \code{sudo systemctl daemon-reload \&\& sudo systemctl restart leo}.
\end{enumerate}

\noindent \textbf{Verify both halves:}
\begin{lstlisting}[language=bash]
ssh pi@$ROBOT_HOST 'chrt -p $(pgrep -f serial_node | head -1)'
\end{lstlisting}
\textbf{Expected:} \code{scheduling policy: SCHED\_FIFO},
\code{priority: 25}. \\
\textbf{And the payoff:} \code{rostopic hz /firmware/wheel\_odom} should read
$\sim19$\,Hz. Before the fix it read $1$\,Hz.

\subsection{Change 2: camera settings that reset on every restart}
\label{ap:repro-camreset}
\emph{Why:} the D455's infrared projector paints a dot pattern that helps
depth and \emph{ruins} the checkerboard and LED detection. It is therefore
turned off, and the depth compression level is lowered to save Pi CPU. Both
are \code{dynamic\_reconfigure} values, so \textbf{they revert to factory
defaults every time the camera node restarts}, and the camera node restarts
whenever the stack does.

\begin{lstlisting}[language=bash]
rosrun dynamic_reconfigure dynparam set /camera/stereo_module laser_power 0
rosrun dynamic_reconfigure dynparam set \
  /camera/depth/image_rect_raw/compressedDepth png_level 1
\end{lstlisting}
\code{tools/restart\_stack.sh} and the watchdog both reapply these, which is
why you rarely notice. If your checkerboard detection suddenly degrades, check
\code{laser\_power} first:
\begin{lstlisting}[language=bash]
rosrun dynamic_reconfigure dynparam get /camera/stereo_module | grep laser_power
\end{lstlisting}

\subsection{Change 3: disable the robot's own rosbridge}
\label{ap:repro-phantom}
\emph{Why:} the stock image runs a \code{rosbridge\_server} on the robot for
its factory web UI. This project's cockpit connects to a rosbridge on the
\emph{PC} instead, so the robot's copy has zero clients, yet it was measured
consuming $19\%$ of a CPU core on a Pi already thermally throttled. Its
consumption \emph{grows with uptime} (a leak: a freshly respawned instance sits
at $0\%$), and its \code{unregister\_timeout} of $86400$\,s let it live outside
the ROS graph, invisible to \code{rosnode list}.

Comment the node out in \code{leo\_bringup.launch} on the robot:
\begin{lstlisting}[language=xml]
<!-- DISABLED: unused, and leaks CPU on a thermally limited Pi.
<node name="rosbridge_server" pkg="rosbridge_server"
      type="rosbridge_websocket">
  <param name="unregister_timeout" value="86400"/>
</node>
-->
\end{lstlisting}
Then \code{sudo systemctl restart leo}. Killing the process without editing
the launch file is futile: \code{rosmon} respawns it within seconds.

\section{The code we wrote, node by node}
\label{ap:repro-ourcode}
Everything upstream (openVINS, MINS, the camera driver) is unmodified. The
project-specific logic lives in six files, all in
\code{catkin\_ws/src/leo\_navigation/scripts/} except the last. Each has a long
header comment explaining \emph{why} it exists: read that header before
changing anything, because most of these nodes exist to work around a specific
measured defect, and the workaround looks arbitrary until you know which.

\subsection{imu\_sanitizer.py: the IMU is not trustworthy raw}
\textbf{Subscribes} \code{/firmware/imu} (\code{leo\_msgs/Imu}) \quad
\textbf{Publishes} \code{/imu/data\_clean} (\code{sensor\_msgs/Imu})

This node exists because the CORE2 IMU has \emph{four} separate defects, each
found by measurement, and both estimators consume its output.

\begin{enumerate}[leftmargin=1.4em]
\item \textbf{Isolated glitches.} Roughly two samples in nine thousand come out
  absurd (one accelerometer reading of $5\times10^{16}$\,m/s$^2$ was recorded).
  A single such sample destroys any filter instantly. Samples outside
  \code{\textasciitilde gyro\_limit} ($35$\,rad/s) or
  \code{\textasciitilde accel\_limit} ($160$\,m/s$^2$) are replaced by the last
  good one, keeping the cadence intact.
\item \textbf{A frozen-firmware mode.} After a serial desync the board can
  retransmit the \emph{same} sample forever: bit-identical values, which a
  real sensor never produces. $N$ identical samples in a row triggers an
  automatic \code{/firmware/reset\_board} call (throttled).
\item \textbf{Accelerometer scale error.} With the robot verified stationary
  ($502$ samples, wheel odometry $0.000$\,m/s) the raw magnitude read
  $8.7525$\,m/s$^2$ against a true local gravity of $9.790$. That is $-10.6\%$,
  and it is in the sensor: raw and sanitised agreed to $0.005$. Corrected by
  \code{\textasciitilde accel\_scale} $= 9.790/8.7525 = 1.1186$.
\item \textbf{Gyroscope bias, badly out of spec.} Stationary, the gyro read
  $3.679$\,deg/s of total bias where a healthy MEMS part sits at
  $0.01$--$0.1$. Because attitude is the \emph{integral} of the gyro, that
  mis-projects gravity into the body frame and injects a phantom horizontal
  acceleration far larger than the accelerometer error. Subtracted, and
  \textbf{re-measured at every node start} (\code{\textasciitilde gyro\_autocal},
  default \code{true}) because bias drifts with temperature and this Pi runs at
  $86\,^\circ$C.
\end{enumerate}

\noindent \textbf{Parameters you may need:}
\begin{lstlisting}[language=bash]
~accel_scale    1.1186   # 1.0 disables; see (3) above
~gyro_scale     1.0      # NOT measured yet -- see below
~gyro_autocal   true     # re-measure bias at startup if the robot is still
~gyro_limit     35.0     # rad/s, glitch rejection
~accel_limit    160.0    # m/s^2
\end{lstlisting}

\noindent \textbf{The honest gap.} \code{\textasciitilde gyro\_scale} is
\textbf{deliberately left at $1.0$, i.e.\ no correction.} A stationary
calibration measures the \emph{bias} but is mathematically incapable of
measuring the \emph{scale}: at rest the true rate is zero, so the scale factor
drops out of the observation entirely. Measuring it requires a turn of known
amplitude against physical ground truth: see
\S\ref{ap:repro-calibrun}. Setting this parameter on a guess would corrupt
\emph{both} estimators at once.

\subsection{wheel\_remap.py: a two-line bug that costs a day}
\textbf{Subscribes} \code{/joint\_states} \quad
\textbf{Publishes} a 2-element \code{JointState} MINS can read

MINS reads wheel velocities \textbf{by array index, not by joint name}:
\code{data.m1 = msg->velocity.at(0)} with a comment claiming it is the front
left wheel. This rover publishes its four wheels in a different order, so MINS
was silently fed the wrong wheels. Nothing errors; the estimate is simply
wrong. This node republishes a two-element message in the order MINS assumes.

\begin{lstlisting}[language=bash]
~front_left_name    ~front_right_name     # which joints to pick
~rear_left_name     ~rear_right_name
~in_topic  /joint_states                  ~out_topic
\end{lstlisting}

\noindent \textbf{Lesson worth carrying:} when integrating any upstream
estimator, verify how it identifies its inputs. Index-based assumptions are
common and fail silently.

\subsection{pose\_selector.py: one topic that never breaks}
\textbf{Subscribes} \code{/mins/imu/odom}, \code{/ov\_msckf/odomimu} \quad
\textbf{Publishes} \code{/robot\_pose\_fused} \quad
\textbf{Service} \code{/pose\_selector/set\_source} (\code{std\_srvs/SetBool};
\code{true} $=$ MINS)

Downstream consumers (the cockpit, logging, any navigation) need
\emph{one} pose topic that keeps working when you switch estimators. The two
estimators have independent origins, so a naive switch would produce a jump of
tens of metres. This node applies an SE(3) correction at the moment of the
switch so the output is continuous.

\begin{lstlisting}[language=bash]
~default_source        mins    # what is served at startup
~max_plausible_speed           # rejects physically impossible jumps
~mins_topic  ~vins_topic  ~odom_frame  ~base_frame  ~imu_frame
\end{lstlisting}

\noindent \textbf{Important limitation, and it bites:} the SE(3) correction is
computed \emph{on a switch only}. If MINS itself respawns it re-initialises at
a new origin, and that jump is \textbf{not} absorbed: the fused topic jumps
with it. This is one reason the estimators are deliberately \emph{not} allowed
to bring down the stack: driving must never depend on estimator health.

\subsection{carolus\_vicon\_bridge.py: the only absolute reference}
\textbf{Publishes} \code{geometry\_msgs/PoseStamped} on the topic MINS accepts
as a VICON input

Every other measurement on this platform is relative. The Carolus beacon, a
fixed LED constellation at a known place, gives a drift-free global fix
whenever it is in view. This node reformats that fix into the type MINS's
external-reference slot consumes, making it an ``indoor GPS'' that cancels
unbounded drift.

\begin{lstlisting}[language=bash]
~global_frame   ~body_frame   ~out_topic   ~rate_hz   ~max_tf_age
\end{lstlisting}

\noindent \textbf{Rule established 2026-07-08 and violated once since:} the
topic \code{/mins/external\_ref/carolus} is the \textbf{private slot} of this
bridge. Nothing else may publish on it. A trajectory recorder that also wrote
there produced a genuinely confusing failure.

\subsection{leo\_backend.py: the cockpit's brain}
This is the largest file in the project ($\sim4500$ lines) and it runs on the
workstation, launched by \code{leo\_backend\_launch.sh}. It is \emph{not} a
normal ROS node in spirit: it is the bridge between the browser and ROS.

\begin{itemize}[leftmargin=1.4em]
\item \textbf{Inbound:} JSON commands on \code{/mission/command}. Every cockpit
  button publishes one, e.g.\ \code{\{"action":"set\_pose\_source",
  "source":"VINS"\}}. To add a button, add an \code{elif action == "..."}
  branch in \code{\_on\_command} and a click handler in the HTML.
\item \textbf{Outbound:} a JSON telemetry blob at $10$\,Hz on
  \code{/mission/telemetry}, plus human-readable lines on
  \code{/mission/log}. The browser subscribes to both through rosbridge.
\item \textbf{Also owns:} the autonomous patrol state machine, the beacon
  approach logic, the trajectory buffers, the detached rosbag recorder, the
  MATLAB export, and the ground-truth calibration panel.
\end{itemize}

\noindent \textbf{Two hard-won implementation notes.} Launching a GUI
application (MATLAB) from this process needs three separate things or it dies
within seconds: a double fork so it survives its parent, a pseudo-terminal
(MATLAB's \code{-r} exits when stdin is not a tty), and a reconstructed
\code{XDG\_RUNTIME\_DIR}/\code{DBUS\_SESSION\_BUS\_ADDRESS} because the backend
is started without a login session. And any recording that matters is written
by a \emph{detached} \code{rosbag} process, because this backend's own buffers
are lost if it restarts.

\subsection{vins\_rotation\_guard.py: deliberately not used}
\label{ap:repro-guardfile}
This file is still in the tree and is \textbf{intentionally not launched.} It
was written to freeze openVINS's reported position during turns, on the belief
that openVINS diverged when rotating. \textbf{It made things measurably worse}
(\S\ref{ap:repro-guard}). It is kept as a record of the reasoning error, not as
a component. Do not re-enable it without reading that section.

\section{Configuration files, and the traps inside them}
\label{ap:repro-config}
\subsection{A trap that applies to every YAML here}
\label{ap:repro-yamltrap}
\textbf{Read this before editing any estimator config.} Both openVINS and MINS
parse their YAML with OpenCV's \code{FileStorage} (the files begin with
\code{\%YAML:1.0}), which is \emph{not} a general YAML parser. It
\textbf{aborts on a long comment placed after a value on the same line}.

The failure mode is vicious: the node starts, subscribes to nothing, and
respawns forever, with no error message naming the file. This trap was hit
twice on this project: the second time \emph{after} it had already been
documented.

\begin{itemize}[leftmargin=1.4em]
\item Keep inline comments \textbf{short}, or put them on their own line.
\item Validate before restarting anything:
\begin{lstlisting}[language=bash]
cd ~/TOUT/catkin_ws/src/open_vins/config/leo
python3 -c "import cv2; f=cv2.FileStorage('estimator_config.yaml', \
cv2.FILE_STORAGE_READ); print('readable:', f.isOpened()); f.release()"
\end{lstlisting}
  \textbf{Expected:} \code{readable: True}.
\item \textbf{Never test a config from \code{/tmp}}: relative paths inside it
  resolve differently and you get a false negative.
\end{itemize}

\subsection{openVINS}
Three files, all in \code{catkin\_ws/src/open\_vins/config/leo/}:
\begin{lstlisting}[language=bash]
estimator_config.yaml     # the filter: init, ZUPT, features, noise, topics
kalibr_imucam_chain.yaml  # camera intrinsics + camera-IMU extrinsics
kalibr_imu_chain.yaml     # IMU noise densities and random walks
\end{lstlisting}

\noindent \textbf{Launch footgun.} The config path is passed as
\code{args="..."}, i.e.\ \code{argv[1]}, and \emph{not} as a
\code{<param name="config\_path">}. On the ROS1 build, the parameter read is
inside an \code{\#if ROS\_AVAILABLE == 2} block and never executes. Using
\code{<param>} fails silently: \code{argv[1]} becomes roslaunch's own
\code{\_\_name:=...} argument, and openVINS reports it cannot open a file with
that literal name.

\noindent \textbf{The values that matter on this platform}, each with the
measurement behind it:

\begin{table}[htbp]
\centering\small
\caption{Key openVINS settings and why they are set this way.}
\label{tab:appovcfg}
\begin{tabular}{@{}llp{6.4cm}@{}}
\toprule
Key & Value & Why \\
\midrule
\code{try\_zupt} & \code{true} &
  \textbf{The single largest fix.} With \code{false} a run diverged to
  $41\,824$\,m; with \code{true}, $0.002$\,m. Zero-velocity updates anchor a
  ground vehicle that stops often. \\
\code{gravity\_mag} & \code{9.790} &
  True local gravity (Melbourne FL). Leaving $9.81$ against a mis-scaled
  accelerometer compounds the error. \\
\code{num\_opencv\_threads} & \code{1} &
  A mutex crash fired $1324$ times, once every $46$\,s. Root cause: OpenCV
  threading. Zero crashes in $7$ minutes after. \\
\code{max\_slam} & \code{0} &
  Pivoting in place makes landmark depth unobservable; matches the MINS
  setting. \\
\code{calib\_cam\_extrinsics} & \code{false} &
  Locked $2026$-$07$-$29$. Online refinement of a geometry that is poorly
  observable under planar motion lets it wander. \\
\code{calib\_cam\_timeoffset} & \code{false} &
  Same reasoning. \\
\code{calib\_cam\_intrinsics} & \code{false} &
  Frozen from Kalibr; offline measurement beats online refinement. \\
\code{init\_imu\_thresh} & \code{0.4} &
  The jerk needed to initialise. The rover must be nudged
  (\S\ref{ap:repro-launch}). \\
\code{use\_stereo} / \code{max\_cameras} & \code{true} / \code{2} &
  Both D455 infrared cameras. \\
\bottomrule
\end{tabular}
\end{table}

\noindent The IMU noise densities in \code{kalibr\_imu\_chain.yaml}
(\code{accelerometer\_noise\_density: 5.0e-03},
\code{gyroscope\_noise\_density: 3.4e-04}) are \emph{measured} values inflated
by a safety factor of three, not datasheet figures.

\subsection{MINS}
A per-sensor tree selected by \code{mins.launch}'s \code{config} argument
(default \code{leo}), under \code{mins\_ws/src/MINS/mins/config/leo/}:
\begin{lstlisting}[language=bash]
config.yaml            # top level: includes the files below
config_imu.yaml        # IMU noise model
config_camera.yaml     # camera intrinsics + extrinsics
config_wheel.yaml      # wheel model: effective radius, wheel base
config_estimator.yaml  # filter and update options
config_init.yaml       # static IMU + wheel initialisation
config_system.yaml     # frames, topics, the VICON slot used by Carolus
\end{lstlisting}
MINS initialises \textbf{standing still} (static IMU + wheel init, $4$--$12$\,s),
which is why it is the default served source. Its transforms are deliberately
isolated onto \code{/tf\_mins} so it cannot fight openVINS over the main
\code{/tf} tree.

\noindent \textbf{Convention trap.} Kalibr and openVINS express the
camera--IMU extrinsic as $\mathbf{T}_{\text{cam,imu}}$; MINS wants the
\textbf{inverse}, $\mathbf{T}_{\text{imu,cam}}$. Copying one into the other
without inverting produces a plausible-looking, wrong geometry. The helper
\code{scripts/fill\_mins\_camchain.py} does the conversion.

\subsection{Where each estimator keeps its configuration}
\label{ap:repro-config-ref}
Each estimator reads a small tree of YAML files; the launch files select
which tree by name.

\paragraph{openVINS.} A single estimator config plus its Kalibr chains, all
under \code{catkin\_ws/src/open\_vins/config/leo/}:
\begin{lstlisting}[language=bash]
config/leo/estimator_config.yaml    # filter: init, ZUPT, feature counts,
                                    #   noise, do_calib_* flags, topics
config/leo/kalibr_imucam_chain.yaml # camera intrinsics + T_imu_cam extrinsics
config/leo/kalibr_imu_chain.yaml    # IMU noise densities / random walks
\end{lstlisting}
The launch node \code{ov\_msckf/run\_subscribe\_msckf} takes the config path
as \code{argv[1]} (a launch \code{<param name="config\_path">} is read only on
the ROS2 build, a real footgun documented in
\code{navigation\_supervision.launch}). Key platform-specific values
(\S\ref{sec:ovinit}, \S\ref{sec:extrinsic}): \code{init\_dyn\_use: false} and
\code{init\_imu\_thresh: 0.4} (jerk-triggered initialisation, the rover
must be nudged to initialise VINS), \code{use\_stereo: true},
\code{max\_cameras: 2}, and the frozen intrinsics
(\code{do\_calib\_int: false}) with online extrinsic/time-offset refinement
(\code{do\_calib\_ext: true}, \code{do\_calib\_dt: true}).

\paragraph{MINS.} A per-sensor tree selected by the \code{config} launch arg
(\code{mins.launch}: \code{<arg name="config" default="leo"/>}), under
\code{mins\_ws/src/MINS/mins/config/leo/}:
\begin{lstlisting}[language=bash]
config/leo/config.yaml           # top-level: includes the files below
config/leo/config_imu.yaml       # IMU noise model
config/leo/config_camera.yaml    # camera intrinsics + extrinsics
config/leo/config_wheel.yaml     # differential-drive wheel model (radius, base)
config/leo/config_estimator.yaml # filter / update options
config/leo/config_init.yaml      # static IMU+wheel initialisation
config/leo/config_system.yaml    # frames, topics, VICON slot (Carolus)
\end{lstlisting}
MINS initialises \emph{standing still} (static IMU+wheel init, 4--12\,s,
\S\ref{sec:ovinit}), which is why it is the default served source. Its
transform output is isolated to \code{/tf\_mins}
(\code{mins.launch}, \S\ref{sec:chain}); the Carolus global fix enters through
the VICON slot on \code{/mins/external\_ref/carolus}.

\paragraph{Beacon / AprilTag.} The autonomy beacon config lives with the
project package: \code{catkin\_ws/src/leo\_autonomy/config/apriltag.yaml}
(family \code{tag36h11}, full-resolution \code{tag\_decimate: 1.0}) and
\code{tags.yaml} (the physical tag: \code{id: 285}, \code{size: 0.2032}\,m).
The Carolus blob detector's parameters are in
\code{carolus\_ws/src/launch/carolusBlobDetection.launch}, including the four
LED \code{known\_points} (\S\ref{sec:july27carolus}); the cockpit's Carolus
panel (\S\ref{sec:july27export}) parses this same file.

\section{Launch file architecture}
\label{ap:repro-launcharch}
\emph{Why this section exists:} the configuration files tell each program what
to believe; the launch files decide \emph{which programs exist, in what order,
and what happens when one of them dies}. That last part is a design decision on
this platform, not a default, and it was arrived at by watching the stack fail.

Everything on the workstation starts from one file. Figure~\ref{fig:applaunch}
shows the inclusion tree.

\begin{figure}[htbp]
\centering
\begin{tikzpicture}[
  font=\small, x=1cm, y=1cm,
  lf/.style={draw, rounded corners=2pt, align=left, inner sep=4pt,
             font=\scriptsize\ttfamily, fill=blue!7},
  grp/.style={draw, rounded corners=2pt, align=left, inner sep=4pt,
              font=\scriptsize, fill=gray!6},
  nd/.style={draw, rounded corners=2pt, align=left, inner sep=3pt,
             font=\scriptsize\ttfamily, fill=green!9},
  edge/.style={-{Latex[length=2mm]}, semithick, gray!70}
]
  % --- entry point
  \node[lf, fill=orange!14] (master) at (0,0)
    {navigation\_master.launch\\[1pt]\normalfont\scriptsize\itshape the only entry point};

  % --- three includes
  \node[lf] (tf)   at (5.6, 1.5) {tf\_static\_extrinsics.launch};
  \node[lf] (mins) at (5.6, 0.0) {mins.launch};
  \node[lf] (sup)  at (5.6,-1.6) {navigation\_supervision.launch};

  \draw[edge] (master.east) -- (tf.west);
  \draw[edge] (master.east) -- (mins.west);
  \draw[edge] (master.east) -- (sup.west);

  % --- what each owns
  \node[nd] (tfn)  at (11.0, 1.5) {static TF publishers\\(base\_link $\to$ camera\_link)};
  \node[nd, fill=yellow!14] (minsn) at (11.0, 0.0)
    {mins\_subscribe\\\normalfont\scriptsize respawn, \code{nice -n -10}};
  \node[nd] (supn) at (11.0,-1.9)
    {imu\_sanitizer · wheel\_remap\\carolus\_vicon\_bridge · ov\_msckf\\
     pose\_selector · leo\_backend\\web\_video\_server · republish$\times$4\\
     camera\_info relay$\times$4 · throttle};

  \draw[edge] (tf.east)   -- (tfn.west);
  \draw[edge] (mins.east) -- (minsn.west);
  \draw[edge] (sup.east)  -- (supn.west);

  \node[font=\scriptsize, align=left, draw, dashed, rounded corners,
        inner sep=4pt, fill=gray!4] at (5.6,-3.9)
    {\textbf{Every node above is} \code{respawn="true"} \textbf{with}
     \code{respawn\_delay="5"}.\\
     None is \code{required="true"} --- see \S\ref{ap:repro-respawn} for why
     that changed.};
\end{tikzpicture}
\caption{The launch inclusion tree on the workstation. One entry point,
three includes, and a deliberate absence: no node is allowed to take the stack
down with it. The robot's own nodes are \emph{not} here: they come from
\code{leo\_bringup.launch} running on the Pi
(\S\ref{ap:repro-robot}).}
\label{fig:applaunch}
\end{figure}

\subsection{What each file owns}
\begin{itemize}[leftmargin=1.4em]
\item \code{navigation\_master.launch}: the only file you should ever launch
  directly. It declares one argument, \code{config} (default \code{leo}), which
  it forwards to \code{mins.launch}, and then includes the three files below. It
  contains \emph{no} nodes of its own.
\item \code{tf\_static\_extrinsics.launch}: the static coordinate transforms
  the platform would otherwise be missing, principally
  \code{base\_link} $\rightarrow$ \code{camera\_link}. Separated so the geometry
  can be corrected without touching the estimators.
\item \code{mins.launch}: the MINS estimator, and nothing else.
\item \code{navigation\_supervision.launch}: everything else: the four image
  \code{republish} nodes and four \code{camera\_info} relays that bring the
  robot's compressed streams into a local namespace, our four Python nodes,
  \code{ov\_msckf}, the cockpit backend and the video server.
\end{itemize}

\noindent \textbf{It does not start \code{robot\_state\_publisher}}, and that is
correct: the URDF and the robot's own drivers belong to
\code{leo\_bringup.launch}, which runs on the Pi. Every file in this package
runs on the workstation.

\subsection{Reliability: why nothing is \code{required="true"}}
\label{ap:repro-respawn}
\textbf{This is the most important design decision in the launch layer, and it
is the reverse of what it was.}

\code{roslaunch} offers two policies when a node dies.
\code{required="true"} means \emph{if this node exits, kill everything};
\code{respawn="true"} means \emph{restart it and leave the rest alone}.

MINS was originally \code{required="true"}, on the reasonable-sounding argument
that restarting an estimator mid-flight without re-initialisation is
meaningless, so a crash should be loud. On $2026$-$07$-$06$ that policy was
observed doing real damage: MINS aborted on an assertion
(\code{SystemManager.cpp:395}, a backward IMU timestamp), and because it was
\code{required}, roslaunch tore down the entire stack, including
\code{leo\_backend}. The operator lost the ability to \emph{drive the robot}
because an estimator had hiccupped.

The principle established from that, and worth carrying to any similar system:

\begin{quote}
\textbf{Driving must never depend on estimator health.}
\end{quote}

Every node is now \code{respawn="true"} with \code{respawn\_delay="5"}. The
five-second delay matters: without it, a node that fails at startup (a bad
config, a missing topic) respawns in a tight loop that saturates a CPU and
floods the log, which on a thermally limited platform is its own failure.

\noindent \textbf{The cost, stated plainly.} On respawn, MINS re-initialises at
a \emph{new origin}, so \code{/mins/imu/odom} jumps. The
\code{pose\_selector}'s SE(3) correction does \textbf{not} absorb that jump:
it only compensates at the moment of a deliberate source switch. So a
mid-run MINS respawn produces a visible discontinuity in
\code{/robot\_pose\_fused}. That is an accepted trade-off: a discontinuous
trajectory is recoverable, a dead cockpit during a manoeuvre is not.

\noindent \textbf{A stale comment to be aware of.} The header of
\code{mins.launch} still claims MINS ``stays \code{required="true"}''. It does
not: the node declaration twenty lines below says \code{respawn="true"}, and
a second comment on the same node explains the change. The header was simply not
updated. Trust the declaration, not the prose above it; this is exactly the kind
of contradiction that misleads a successor.

\subsection{Passing a config path: \code{args}, never \code{<param>}}
\label{ap:repro-argsparam}
Both estimators take their YAML path as a \textbf{positional command-line
argument}, and getting this wrong produces one of the most misleading failures
in the whole system.

\begin{lstlisting}[language=xml]
<!-- CORRECT -->
<node name="mins_subscribe" pkg="mins" type="subscribe"
      respawn="true" respawn_delay="5"
      launch-prefix="nice -n -10 stdbuf -oL"
      args="$(find mins)/config/$(arg config)/config.yaml">
</node>

<!-- WRONG: silently ignored on the ROS1 build -->
<node name="mins_subscribe" pkg="mins" type="subscribe">
  <param name="config_path" value="..."/>
</node>
\end{lstlisting}

\noindent \textbf{Why the second form fails, and fails confusingly.} In the
upstream source, the call that would read a \code{config\_path} ROS parameter
sits inside an \code{\#if ROS\_AVAILABLE == 2} block: it is compiled out of
the ROS1 build entirely. The program therefore falls back to reading
\code{argv[1]}. But roslaunch appends its own arguments to every node, so
\code{argv[1]} becomes \code{\_\_name:=mins\_subscribe}, and the estimator
reports that it cannot open a configuration file with that literal name. The
error message names a string you never typed, which sends you looking in the
wrong place.

\noindent \code{ov\_msckf} has exactly the same contract, for exactly the same
reason. Both are documented in the launch files themselves so the trap is not
rediscovered.

\subsection{Two launch-prefix tricks used here}
\begin{itemize}[leftmargin=1.4em]
\item \code{nice -n -10} on MINS. The workstation's own IDE indexers were
  measured starving the estimator of CPU, which produced episodes that looked
  like estimator bugs. A negative nice value requires permission, granted to
  this user in \code{limits.conf}.
\item \code{stdbuf -oL} on MINS. Its initialisation diagnostics go to plain
  \code{stdout}. When roslaunch redirects that to a file, glibc switches from
  line-buffered to fully-buffered, and at this output volume the \code{[INIT]}
  messages can sit unflushed for \emph{minutes}: looking exactly like a
  hung estimator when it is initialising normally. Forcing line buffering makes
  the diagnostics appear as they are printed.
\end{itemize}

\noindent The robot side uses the same idea for a different purpose:
\code{launch-prefix="chrt -f 25"} on \code{serial\_node}, discussed in
\S\ref{ap:repro-rtfix}.

\subsection{How to add a node}
\begin{enumerate}[leftmargin=1.4em]
\item Put the script in
  \code{catkin\_ws/src/leo\_navigation/scripts/} and make it executable
  (\code{chmod +x}). A non-executable script produces
  \code{cannot locate node of type [...]}, which reads like a missing file.
\item Declare it in \code{navigation\_supervision.launch}, with
  \code{respawn="true" respawn\_delay="5"} unless you have a specific reason
  otherwise, and write that reason in a comment.
\item Add \code{output="screen"} if you want its log in
  \code{logs/navigation\_master.log}. Without it, diagnostics are difficult to
  find.
\item Restart with \code{bash tools/restart\_stack.sh}, never a bare
  \code{roslaunch} (\S\ref{ap:repro-sourcetrap}).
\item Verify it registered \emph{and} publishes:
\begin{lstlisting}[language=bash]
rosnode list | grep my_node
rostopic hz /my_topic
\end{lstlisting}
  A node can register and stay silent; only the second command proves it works.
\end{enumerate}

\subsection{Two files the diagram does not show, and why}
\label{ap:launch-notshown}

\paragraph{\code{pose\_selector.launch}.} The diagram above shows
\code{pose\_selector} as a node inside
\code{navigation\_supervision.launch}, which is how it runs in normal
operation. A standalone \code{pose\_selector.launch} also exists, for
bringing the selector up alone against an already-running stack when you are
debugging source switching. Whichever route you use, one parameter decides
whether the TF tree is well formed:
\begin{lstlisting}[language=XML]
<param name="base_frame" value="base_footprint"/>
\end{lstlisting}
Set to \code{base\_link} instead, it gives that frame a second parent and the
pose starts alternating silently between two coherent answers
(\S\ref{ap:trap-tf}). Check it in both files if you ever edit one.

\paragraph{\code{carolus\_tf\_bridge.py}.} Absent from every launch file, and
deliberately so. It is the beacon-anchored absolute-error bridge of
\S\ref{ap:repro-absolute}, and it freezes a static
\code{beacon\_link}~$\to$~\code{odom} transform at the first fix. A node that
silently re-anchors the world frame has no business starting automatically
alongside an unrelated run, so it is started by hand, by an operator who
intends to use it:
\begin{lstlisting}[language=bash]
rosrun leo_navigation carolus_tf_bridge.py
\end{lstlisting}
If you later decide to add it to a launch file, leave
\code{\textasciitilde publish\_body\_chain} at its default of false: enabling
it re-parents \code{base\_link} and corrupts the tree in the same way
described above.

\section{Launching the system, step by step}
\label{ap:repro-launch}
\emph{Why an order:} the master must exist before clients, and the estimators
must see sensor data from their first sample or they initialise on nothing. The
\code{leo} command sequences all of this; the value of knowing the order is
being able to diagnose a \emph{partial} start, which is the common case.

\subsection{The nominal sequence}
\begin{enumerate}[leftmargin=1.4em]
\item \textbf{Power the robot} and wait $\sim40$\,s. Its own \code{leo.service}
  brings up the ROS master and the sensors before anything on the PC can
  connect.
\item \textbf{Check the link first.} From any directory:
\begin{lstlisting}[language=bash]
cd ~/TOUT && source tools/robot_env.sh
ping -c 3 "$ROBOT_HOST"
\end{lstlisting}
  \textbf{Expected:} \code{0\% packet loss}. If this fails, stop: everything
  after it will fail in a more confusing way.
\item \textbf{Start the workstation side.} Two scripts, in this order. The
  first brings up the web stack (static server on $8000$, rosbridge on
  $9090$, the TLS bridge on $9443$, the video server on $8080$, and
  \code{cloudflared} if configured); the second brings up the navigation
  stack (MINS, openVINS, \code{pose\_selector}, \code{leo\_backend}). Both
  are idempotent, so re-running them is safe.
\begin{lstlisting}[language=bash]
cd ~/TOUT && source tools/robot_env.sh
touch /tmp/leo_maintenance          # suspend the watchdog while starting
bash web/start_web.sh
bash tools/restart_stack.sh
rm -f /tmp/leo_maintenance          # DO NOT FORGET THIS
\end{lstlisting}
  \textbf{Expected:} a series of \code{[ok]} lines from the first script and
  a launch log from the second. No \code{ECHEC} line.

  \emph{Why the maintenance flag:} the cron watchdog restarts anything it
  believes is dead, and a stack that is halfway up looks exactly like a stack
  that has died. Without the flag the two of you will fight
  (\S\ref{ap:repro-watchdog}). The flag is a plain file, so \textbf{if you
  forget to remove it the watchdog stays disabled indefinitely} and you lose
  every automatic recovery in this document.

\item \textbf{Confirm the robot is seen:}
\begin{lstlisting}[language=bash]
ping -c 1 -W 1 "$ROBOT_HOST" && echo "robot online"
ss -ltn | grep -E ':(8000|8080|9090|9443) '   # four listening ports
\end{lstlisting}
  \textbf{Expected:} \code{robot online}, and four ports bound.

  \emph{A note on \code{leo start}:} the workstation used for this project
  carried a convenience wrapper at \code{\textasciitilde/bin/leo} that
  performed the sequence above plus a status view. It lives outside the
  project tree and is \textbf{not} part of the repository, so it will not
  exist on your machine. The commands above are what it actually ran; use
  them directly, or write your own wrapper around them.
\item \textbf{Confirm the four nodes that matter registered:}
\begin{lstlisting}[language=bash]
source /opt/ros/noetic/setup.bash
rosnode list | grep -E "mins_subscribe|ov_msckf|pose_selector|imu_sanitizer"
\end{lstlisting}
  \textbf{Expected:} all four. A missing \code{ov\_msckf} almost always means
  it aborted at startup: look in \code{logs/navigation\_master.log}, and
  suspect the YAML trap of \S\ref{ap:repro-yamltrap} first.
\item \textbf{Check the \emph{rates}, not just the names.} A registered node
  can be completely silent; this is the real health test:
\begin{lstlisting}[language=bash]
rostopic hz /imu/data_clean                # expect ~80 Hz
rostopic hz /mins/imu/odom                 # expect ~80-90 Hz
rostopic hz /wheel_odom_with_covariance    # expect ~19-20 Hz
rostopic hz /pc/camera/infra1/image_rect_raw   # expect ~15 Hz
\end{lstlisting}
  \code{/wheel\_odom\_with\_covariance} at $1$\,Hz instead of $19$ is the
  signature of the UART problem of \S\ref{ap:repro-uart}.
\item \textbf{Open the cockpit} at \url{http://localhost:8000/ops.html}.
\end{enumerate}

\subsection{Making openVINS start: the classic beginner trap}
MINS initialises standing still. \textbf{openVINS does not.} It waits for an
accelerometer jerk above \code{init\_imu\_thresh} before publishing anything
\emph{at all}. Newcomers conclude openVINS is broken when it is merely waiting.

\begin{enumerate}[leftmargin=1.4em]
\item Drive the rover briskly forward, then back, about one second each way.
\item Confirm it woke up:
\begin{lstlisting}[language=bash]
rostopic hz /ov_msckf/odomimu           # expect ~75-90 Hz
\end{lstlisting}
  Silence means it has not initialised. Jerk it again: more sharply.
\end{enumerate}

\subsection{Restarting cleanly when something is wrong}
Do \textbf{not} kill individual nodes as a first response; several are
respawn-supervised and will come back in a state you did not intend.
\begin{lstlisting}[language=bash]
cd ~/TOUT
bash tools/restart_stack.sh
\end{lstlisting}
\textbf{Expected:} \code{stack relancée}, then
\code{MINS initialisé après <N>s}, then \code{switched to MINS}, then
\code{TERMINE}. Typical $N$ is $8$--$12$\,s.

This script also raises the maintenance flag while it works, waits for MINS to
publish before declaring success, and reapplies the camera settings that reset
on restart (\S\ref{ap:repro-camreset}).

\subsection{A short tour of the cockpit}
\begin{itemize}[leftmargin=1.4em]
\item \textbf{Operations}: live video, mission state, battery, and the
  \emph{Pose Source} switch (MINS/VINS). Manual driving is here.
\item \textbf{Trajectory}: the live plot of the served pose, the
  \emph{Test 1 / Test 2} selector, the MATLAB export buttons, and the
  \emph{ground-truth calibration} panel (\S\ref{ap:repro-calibrun}).
\item \textbf{Logbook}, the project journal and the download card for this
  report.
\item \textbf{PID / Nav Modes / Robot}, tuning, autonomy state, and the URDF
  view.
\end{itemize}
The cockpit is served over an authenticated Cloudflare tunnel as well, so it
works from a phone or tablet standing next to the robot, which is how the
calibration is meant to be done.

\subsection{Editing the cockpit without breaking it}
\label{ap:repro-editing-web}
Two habits, both learned the hard way, and both cheap.

\paragraph{Your edit is live but you are not seeing it.} This is the first
thing that will happen to you, and it is not your edit. Cloudflare caches
aggressively and a dashboard rule overrides the \code{Cache-Control:
no-cache} that \code{serve.py} sends. Static assets therefore carry a version
stamp in their query string:
\begin{lstlisting}[language=html]
<script src="app.js?v=r37"></script>
\end{lstlisting}
The query string is part of Cloudflare's default cache key, so
\textbf{incrementing that number is what forces a fresh copy}. Bump it in
every page that references the file you changed, and keep the stamps
consistent across pages, since they have drifted apart before and produced
two pages running different versions of the same script.

A browser force-reload does \emph{not} help, because the stale copy sits at
the network edge and not on your machine. If you are testing over
\code{http://localhost:8000} you will not see the problem at all, which is
exactly why it surprises people the first time they check the public URL.

\paragraph{Validate JavaScript before you reload.} \code{app.js} is large and
a brace inserted in the wrong place produces a file that loads without error
and simply stops executing partway through, so the page renders and half the
panels are dead. Counting braces by hand does not catch this; it failed once
in this project and cost an afternoon. Parse the file instead:
\begin{lstlisting}[language=bash]
pip install esprima                     # once
python3 -c "import esprima,sys; esprima.parseScript(open('web/app.js').read()); print('OK')"
\end{lstlisting}
Do the same for the inline scripts in \code{ops.html} if you touch them. A
parse takes under a second and turns a silent partial failure into a line
number.

\subsection{Two pages, deliberately different: control versus demonstration}
\label{ap:repro-twopages}
\noindent\textbf{Two pages, deliberately different (cockpit v4).}
\code{web/ops.html} is the control surface: it holds every command path,
manual driving, AUTO engagement, pose-source switching, beacon reset, MATLAB
export. \code{web/demo.html} subscribes to the same topics and publishes
nothing at all. It has no command bindings to disable, because it has none to
begin with.

The separation is not cosmetic. A defence is the one occasion where the
machine is driven by someone who is talking to a room rather than watching the
robot, and where a mis-click is expensive and public. \code{demo.html} is
built so that no mis-click exists: it shows the trajectory in the beacon
frame, heading quivers, and the RMS error at $56$\,px, and it cannot move the
rover. Open it on the projector; keep \code{ops.html} on your own screen.

\paragraph{Why read-only is a structural property, not a setting.}
The distinction matters because the two obvious alternatives are both worse.
Disabling the controls with a flag leaves the command paths present in the
page, one \code{localStorage} key or one console call away from being live
again, and it invites the reader of the code to wonder which state the page is
in. Hiding them with CSS is worse still, since the handlers remain bound and a
keyboard shortcut or a stray focus event can fire them. \code{demo.html}
instead \emph{subscribes} and never \emph{advertises}: it holds no publisher,
no command topic and no \code{data-cmd} binding anywhere in its markup. There
is no state in which it can move the robot, so there is no state to verify
before a defence.

\paragraph{What each page is for.}
\begin{center}
\small
\begin{tabular}{@{}lL{4.5cm}L{4.5cm}@{}}
\toprule
 & \code{ops.html} & \code{demo.html} \\
\midrule
Audience     & the operator, alone & a room, during a defence \\
ROS role     & subscriber \emph{and} publisher & subscriber only \\
Commands     & manual drive, AUTO, pose source, beacon reset, MATLAB export & none exist \\
Failure mode & a mis-click moves the rover & a mis-click does nothing \\
Where to open it & your own screen & the projector \\
\bottomrule
\end{tabular}
\end{center}

\noindent\textbf{A practical consequence for the transport.} Because
\code{demo.html} never publishes, it is immune to the entire class of
symptoms described in \S\ref{ap:repro-webtransport}: a zombie TLS port makes
buttons inert, and a page with no buttons cannot present that symptom. If the
demonstration view is updating and the control view is not, you have localised
the fault to the outbound direction of the WebSocket before running a single
command.

\section{Reproducing the experiments}
\label{ap:repro-exp}
\subsection{A MINS versus openVINS comparison run}
\label{ap:repro-comprun}
\emph{Why a rosbag and not the website:} the cockpit's in-memory buffer is lost
if the backend restarts, and rosbridge can desynchronise. A run that matters
must be written to disk by a process that does not depend on the web stack.

\begin{enumerate}[leftmargin=1.4em]
\item \textbf{Pre-flight.} Run the checklist of \S\ref{ap:repro-checklist}.
  This is not optional, most of the discarded data in this project failed
  one of those six items.
\item \textbf{Wake openVINS} with a jerk, and confirm both estimators publish.
\item \textbf{Mark a start point} on the floor with tape. You will drive a
  closed loop and return to it; the gap on return is the only drift measure
  available without ground truth.
\item \textbf{Start recording} (from \code{\textasciitilde/TOUT}):
\begin{lstlisting}[language=bash]
cd ~/TOUT
tools/record_trajectories.sh 180 essai_A
\end{lstlisting}
  Records $180$\,s into
  \code{data/trajectories/essai\_A\_<timestamp>.bag}. Omit the duration to
  record until \code{Ctrl-C}. \textbf{Both} estimators are always captured
  together: that is the entire point.
\item \textbf{Drive the loop} and return to the marked point. Drive slowly and
  smoothly; note the commanded speed.
\item \textbf{Convert to CSV:}
\begin{lstlisting}[language=bash]
python3 tools/bag_to_csv.py data/trajectories/essai_A_<timestamp>.bag
\end{lstlisting}
  \textbf{Expected:} one CSV per source, \code{\_mins.csv}, \code{\_vins.csv},
  and \code{\_carolus.csv} if a beacon was visible. Columns are
  \code{t,x,y,z,qx,qy,qz,qw,yaw\_rad}, so any tool can read them: MATLAB is
  a convenience, not a requirement.
\item \textbf{Plot and get the numbers} (run from \code{tools/} so MATLAB finds
  the script):
\begin{lstlisting}[language=bash]
cd ~/TOUT/tools
matlab -batch "plot_trajectories('../data/trajectories/essai_A_<ts>_mins.csv')"
\end{lstlisting}
  \textbf{Expected:} a printed metric block (divergence, scale, path lengths,
  heading gap) and a saved \code{\_comparison.png} plus eight
  \code{\_p1..\_p8.png} panels.
\item \textbf{Put it in the report:} copy \code{\_comparison.png} to
  \code{report\_latex/figures/traj\_comparison.png} and the panels to
  \code{traj\_p1.png}\dots\code{traj\_p8.png}, then recompile
  (\S\ref{ap:repro-report}).
\end{enumerate}

\noindent \textbf{Two independent drives instead of one?} Use
\code{tools/compare\_tests.m} on two runs labelled Test~1 and Test~2. It
deliberately does \emph{not} compute the Kabsch divergence: that requires
both estimators observing the \emph{same} motion at the \emph{same} instants,
which two separate drives are not. It reports loop-closure error and heading
drift instead, which are meaningful for independent runs.

\subsection{How to read the nine-panel figure}
The panels are described one by one in \S\ref{sec:july29panels}, and the
mathematics behind panels~2--4 in \S\ref{sec:july29align} and
\S\ref{sec:july29extent}. Two habits will keep you honest:

\begin{itemize}[leftmargin=1.4em]
\item \textbf{Read panels 2 and 4 together, never separately.} Panel~2 aligns
  the trajectories \emph{including a scale factor}, so it is constructed to
  hide a scale disagreement. Panel~4 is the only one that exposes it. On the
  July~29 run, $47\%$ of the residual was pure scale error.
\item \textbf{The scale factor was fitted to MINS, and MINS is not ground
  truth.} The script prints a warning whenever the fitted scale departs from
  unity by more than $3\%$. Quoting the scale-corrected residual alone would be
  quoting a number built to be insensitive to the error you are measuring.
\end{itemize}

\subsection{Ground-truth calibration: the most valuable measurement}
\label{ap:repro-calibrun}
Everything above compares estimators to each other. This is the one procedure
that compares them to \emph{reality}, and it is what a new student should do
first.

\begin{enumerate}[leftmargin=1.4em]
\item \textbf{Build the measurement zone.} Follow
  \code{docs/protocole\_calibration\_terrain.md}. The essentials: construct a
  right angle with a 3-4-5 triangle ($90$, $120$, $150$\,cm, a tape measure
  does this to $0.08^\circ$, where a protractor manages $\pm1^\circ$), and hang
  two plumb lines from the robot's centreline so its heading is marked on the
  floor without parallax.
\item \textbf{Open the Trajectory tab} on a phone or tablet and pick a
  manoeuvre: \code{90\textdegree\ G}, \code{90\textdegree\ D}, or \code{1\,m}.
\item \textbf{Mark start}, drive the manoeuvre, \textbf{mark end}. The backend
  freezes that instant: take as long as the measurement needs.
\item \textbf{Measure physically.} For a rotation, read the lateral offset $d$
  of the far plumb mark from the reference line over the lever arm $L$, and
  compute the residual $\varepsilon = \arctan(d/L)$; the true angle is
  $90^\circ \mp \varepsilon$. With $L = 1$\,m and $d$ read to $\pm2$\,mm this
  gives $\pm0.11^\circ$.
\item \textbf{Enter what you measured}, never the value you aimed for. Then
  validate.
\item \textbf{Three passes per manoeuvre, both directions.} The panel reports
  the mean factor and its standard deviation, and refuses an entry that
  disagrees with all four sensors by more than a factor of five.
\end{enumerate}

\noindent \textbf{The reasoning that matters more than the numbers.} The first
campaign found the gyroscope over-reporting by $7\%$ \emph{and} the wheels by
$9\%$. Wheel-derived heading comes from differential odometry and does not use
the gyroscope: these are physically independent instruments. MINS and openVINS
are \emph{not} independent witnesses, since both integrate the same gyro.

Two independent sensors agreeing with each other and disagreeing with the
reference points at \textbf{the reference}. So the conclusion was not ``apply a
$7\%$ correction'' but ``re-measure with a better floor protocol''. Decide
between the three outcomes:

\begin{center}\small
\begin{tabular}{@{}lp{8.2cm}@{}}
\toprule
Result over 3 passes & What it means \\
\midrule
both $\rightarrow 1.00$ & the first campaign was measuring the protractor.
  Change nothing; leave \code{gyro\_scale} at $1.0$. \\
both stay $\approx 0.93$ & a real shared error: applicable, but you must
  then explain why two different physical principles drift by the same amount. \\
they diverge & each has its own defect. Treat separately:
  \code{gyro\_scale} for the gyro, effective radius / wheel base for the
  wheels. \\
\bottomrule
\end{tabular}
\end{center}

\noindent Standard deviation above $5\%$ means the \emph{method} still
dominates: do not apply anything, improve the marking.

\subsection{Protocol T1.1: static drift and variance analysis}
\label{ap:repro-t11}
Chapter~12 (\S\ref{sec:t11}) uses this protocol to separate two things that
look identical in a single trajectory plot but demand opposite fixes:
\emph{noise} (bounded, average it away) and \emph{drift} (unbounded, points
at an unmodelled bias in extrinsics or initialisation). It requires the
rover to do nothing at all, which makes it the cheapest measurement in this
guide and, per Table~\ref{tab:apptools}, a fully automated one.

\begin{enumerate}[leftmargin=1.4em]
\item \textbf{Fresh initialisation.} Restart the stack
  (\S\ref{ap:repro-launch}) rather than reuse a session that has been
  running for a while: T1.1 characterises the estimator from a cold start,
  and a warm one hides exactly the divergence it exists to catch.
\item \textbf{Place the rover and do not touch it again.} Flat ground, clear
  of anything that would tempt you to nudge it once logging starts.
\item \textbf{Log for five minutes} (\code{300}\,s is the script's own
  default, matching the campaign's measurement):
\begin{lstlisting}[language=bash]
cd ~/TOUT/tools
python3 log_pose_static.py 300 /tmp/t1_1_static_pose.csv
\end{lstlisting}
  The console prints \code{Logging /robot\_pose\_fused for 300s -- keep the
  rover PERFECTLY STILL.}, then, after five minutes,
  \code{Wrote <N> samples to /tmp/t1\_1\_static\_pose.csv} with $N$ close to
  $300 \times 20 = 6000$ at the nominal $20$\,Hz publish rate.
\item \textbf{Analyse:}
\begin{lstlisting}[language=bash]
python3 analyze_pose_variance.py /tmp/t1_1_static_pose.csv
\end{lstlisting}
  This reports, per axis: standard deviation (spread around the mean, a
  noise/covariance question if bounded), a least-squares linear drift rate
  in m/s (near zero means noise; a clear nonzero slope means genuine
  unbounded divergence), and a first-10\%-versus-last-10\%-of-the-window
  mean as a second, independent check of the same conclusion.
\end{enumerate}

\noindent \textbf{The verdict this campaign used}: drift below
$1$\,mm/s on every axis passes. Above it, do not touch a noise parameter ---
Eq.~\eqref{eq:chi2}'s $\chi^2$ gate and the covariances upstream of it are
tuned for noise, not for a genuine bias, and loosening them to accommodate an
unbounded drift will make the estimator accept bad data rather than fix the
cause. A failing T1.1 run points at extrinsics or initialisation
(\S\ref{ap:repro-yamltrap}), never at a filter gain.

\subsection{Absolute error: anchoring the trajectory to the beacon}
\label{ap:repro-absolute}

\begin{trapbox}{READ THIS BEFORE PUBLISHING ANY NUMBER FROM THIS SECTION}
\textbf{THE COLOUR CAMERA INTRINSICS ARE WRONG AND MUST BE RE-CALIBRATED
FIRST.} The detector runs with parameters calibrated at $1280\times720$
against a $640\times480$ stream. Its stored principal point,
$c'_x = 643.47$, lies \emph{outside} a $640$-pixel-wide image.

Consequences, derived in \S\ref{sec:july30intrinsics}: every range is
inflated by $69.3\,\%$, and the lateral estimate carries a fixed offset
$\Delta X = \Delta c_x L / w$ that translates the entire scene. The result is
smooth, repeatable, internally consistent and metrically meaningless, which is
the most dangerous combination a measurement can present.

\textbf{The procedure below is correct and ready to run. Its output is not
publishable as a dimensional quantity until the calibration of
\S\ref{ap:next-intrinsics} has been done.} Treat what follows as a rehearsal
until then.
\end{trapbox}

\noindent\textbf{Why this replaces the alignment-based error.} Every figure in
Chapter~12 up to \S\ref{sec:july29align} is a residual left after fitting a
transform between two trajectories, by Kabsch (rigid) or Umeyama (with scale).
\S\ref{sec:july30bridge} shows that such a fit is an orthogonal projection
away from a seven-dimensional subspace $\mathcal{G}$ of translations,
rotations and scale, and that the dominant estimator failure modes, heading
drift and scale drift, are basis vectors \emph{of} that subspace. An
alignment-based metric is therefore blind by construction to the very errors
the comparison exists to expose, and every number it reports is a lower bound.

The bridge removes the fit rather than improving it. The transform is
established once, physically, by anchoring, and thereafter held constant, so
no parameter is estimated at measurement time.

\paragraph{What \code{carolus\_tf\_bridge.py} does.}
It publishes exactly two transforms and derives two topics from them:
\begin{itemize}[leftmargin=1.4em]
\item \code{camera\_frame} $\to$ \code{carolus\_raw}, live, from each optical
  detection of the beacon.
\item \code{beacon\_link} $\to$ \code{odom}, static, captured at the first fix
  and then \textbf{frozen}. This is the anchoring act, and it is the whole
  point: from this instant the odometric world and the optical world share one
  absolute reference.
\item \code{/odom\_in\_beacon} and \code{/carolus\_in\_beacon}, both carrying
  the \emph{robot's} pose in the beacon frame, reached respectively through
  the odometry/estimator chain and through the optical detection. Their
  instantaneous difference is the absolute error.
\end{itemize}

\begin{center}
\begin{tikzpicture}[font=\scriptsize, >=Latex,
  b/.style={draw, rounded corners=2pt, inner sep=3.5pt, align=center}]
  \node[b, fill=accentblue!7, draw=accentblue] (bl) at (0,0) {\code{beacon\_link}\\(frozen anchor)};
  \node[b] (od) at (3.6,0.95) {\code{odom}};
  \node[b] (rob) at (7.4,0.95) {robot pose\\(estimator)};
  \node[b] (cam) at (3.6,-0.95) {\code{camera\_frame}};
  \node[b] (car) at (7.4,-0.95) {robot pose\\(optical)};
  \draw[->] (bl) -- node[above, sloped]{static, captured once} (od);
  \draw[->] (od) -- (rob);
  \draw[->] (bl) -- node[below, sloped]{live detection} (cam);
  \draw[->] (cam) -- (car);
  \draw[<->, alertred, very thick] (rob) -- node[right=2pt, alertred]{$\varepsilon_{\mathrm{abs}}$} (car);
\end{tikzpicture}
\end{center}

\paragraph{Step by step.}
\begin{enumerate}[leftmargin=1.6em]
\item \textbf{Place the beacon} where it stays visible for the whole planned
  loop, and do not move it afterwards. Moving the beacon after anchoring
  invalidates the run silently, because the frozen transform no longer
  describes reality and nothing detects the discrepancy.
\item \textbf{Confirm the detector is actually seeing it} before starting
  anything else:
\begin{lstlisting}[language=bash]
rostopic hz /pose            # non-zero => the beacon is detected
rosnode info /carolus_astrobee_rex | head -20
\end{lstlisting}
  A rate of $0$\,Hz with \code{Not enough blobs < 4} in the log means the LEDs
  are not resolved: adjust distance, exposure or lighting. The bridge cannot
  anchor to something that was never detected.
\item \textbf{Start the bridge by hand.} It is deliberately absent from every
  launch file, so that it can never perturb a run you did not intend it to
  touch:
\begin{lstlisting}[language=bash]
rosrun leo_navigation carolus_tf_bridge.py
\end{lstlisting}
\item \textbf{Verify the anchor took}, and that the TF tree is still
  well-formed (\S\ref{ap:trap-tf}):
\begin{lstlisting}[language=bash]
rostopic echo -n1 /odom_in_beacon
rostopic echo -n1 /carolus_in_beacon
rosrun tf tf_monitor 2>&1 | head -30     # one parent per frame
\end{lstlisting}
\item \textbf{Record both topics}, along with the usual estimator outputs.
  \textbf{This is the step that has never been done}, and its absence is why
  no absolute error appears anywhere in this report:
\begin{lstlisting}[language=bash]
rosbag record -O run_absolute.bag \
  /odom_in_beacon /carolus_in_beacon \
  /mins/imu/odom /ov_msckf/odomimu /robot_pose_fused /firmware/wheel_odom
\end{lstlisting}
\item \textbf{Drive the loop} and return near the start. Keep it short, for
  the thermal reasons of \S\ref{ap:repro-hitrun}.
\item \textbf{Compute the error} as the instantaneous difference between the
  two beacon-frame topics. No alignment step is involved, and none should be
  added: adding one would reintroduce exactly the projection this method
  exists to remove.
\end{enumerate}

\begin{trapbox}{Never let the bridge publish the body chain}
The bridge exposes a \code{\textasciitilde publish\_body\_chain} parameter
which defaults to false and logs an error if enabled. An early revision
published the full body chain in the beacon frame, which gave
\code{base\_link} a second parent and corrupted the tree
(\S\ref{ap:trap-tf}). The lookups do not fail; they alternate between two
plausible answers. Leave the parameter alone.
\end{trapbox}

\subsection{Camera--IMU calibration with Kalibr}
Both estimators consume the same physical calibration, produced offline and
then frozen. Offline measurement beats online refinement: that is why
\code{calib\_cam\_extrinsics} is \code{false}.

\begin{enumerate}[leftmargin=1.4em]
\item \textbf{Target.} Print a Kalibr target (an Aprilgrid is preferable to a
  checkerboard for full-view corner observability) and record its exact
  geometry in a \code{target.yaml}.
\item \textbf{Intrinsics and stereo extrinsics.} Record a bag waving the target
  across the whole field of view of both infrared cameras, then run
  \code{kalibr\_calibrate\_cameras} on the two image topics. This produces the
  \code{*-camchain.yaml}. Achieved on this D455 pair: $0.20$--$0.23$\,px
  reprojection error at a $90.2$\,mm baseline.
\item \textbf{Camera--IMU extrinsics and time offset.} With intrinsics fixed,
  record a second bag \textbf{exciting all IMU axes} and run
  \code{kalibr\_calibrate\_imu\_camera}.
\item \textbf{Install} the outputs: openVINS reads them directly as
  \code{kalibr\_imucam\_chain.yaml} and \code{kalibr\_imu\_chain.yaml}; MINS's
  \code{config\_camera.yaml} is filled from the same camchain via
  \code{fill\_mins\_camchain.py}, remembering the inverse-convention trap.
\item \textbf{Verify} with the live reprojection monitor
  (\code{calibration\_monitor.py}). Acceptance threshold used throughout this
  project: $<0.5$\,px.
\end{enumerate}

\noindent \textbf{Standing caveat, and the recommended first task.} Step~3 has
\textbf{never been completed} on this platform. The deployed camera--IMU
extrinsic is a tape-measure seed. The blocker is physical: the procedure needs
roll and pitch excitation of $\gtrsim1.5$\,rad/s, and a rover on flat ground has
only ever achieved $0.08$. You must \emph{lift and tilt the robot by hand} in
front of the target. Doing this is the single most valuable calibration task
outstanding.

\subsection{Calibration methodology, in full}
\label{ap:repro-calib}
Both estimators consume the same physical calibration, produced offline with
Kalibr~\cite{furgale2013kalibr} and then frozen (offline measurement beats
online refinement, \S\ref{sec:tuning}). The procedure, reproducible from
scratch:

\begin{enumerate}[leftmargin=1.4em]
\item \textbf{Target.} Print a Kalibr calibration target (an Aprilgrid is
  recommended over a checkerboard for full-view corner observability) and
  record its exact geometry in a \code{target.yaml}.
\item \textbf{Intrinsics + stereo extrinsics.} Record a rosbag waving the
  target across the full field of view of both IR cameras, then run
  \code{kalibr\_calibrate\_cameras} on the two image topics. This yields the
  \code{*-camchain.yaml} (intrinsics per camera, plus the stereo baseline).
  On this D455 pair the achieved reprojection error was $0.20$--$0.23$\,px at
  a $90.2$\,mm baseline (\S\ref{sec:tuning}); intrinsics are then frozen
  (\code{do\_calib\_int: false}).
\item \textbf{Camera--IMU extrinsics + time offset.} With the intrinsic
  camchain fixed, record a second bag exciting all IMU axes (rotate and
  translate the rig) and run \code{kalibr\_calibrate\_imu\_camera} with the
  camchain, the \code{imu.yaml} noise model and the target. This produces the
  \code{*-camchain-imucam.yaml} with $^{I}\mathbf{T}_{C}$ and the
  camera--IMU time offset $t_d$. The correct forward-looking mount matters:
  an identity placeholder here was the root cause of the vertical divergence
  fixed in \S\ref{sec:extrinsic} (Eq.~\ref{eq:extr}). $t_d$ and the
  translation are then refined online (\code{do\_calib\_ext},
  \code{do\_calib\_dt}), but the $90^\circ$ rotation must be seeded correctly: it is outside the filter's recoverable region.
\item \textbf{Install.} Copy the Kalibr outputs into each estimator's config:
  openVINS reads them as \code{config/leo/kalibr\_imucam\_chain.yaml} and
  \code{kalibr\_imu\_chain.yaml} directly; MINS's \code{config\_camera.yaml}
  is populated from the same camchain (the prepared
  \code{fill\_mins\_camchain.py} helper maps one to the other). The IMU noise
  densities in \code{kalibr\_imu\_chain.yaml} come from an Allan-variance
  characterisation of the D455 IMU (or the manufacturer figures inflated for
  margin).
\item \textbf{Verify.} The live reprojection-error monitor
  (\code{calibration\_monitor.py}, surfaced on the cockpit's calibration
  panel) reports the current chain's health; the acceptance threshold used
  throughout this project is $<0.5$\,px (\S\ref{sec:tuning}).
\end{enumerate}

\noindent A standing caveat carried in the open-items registry: the deployed
camera--IMU extrinsic is still a tape-measure seed refined online, not a full
\code{kalibr\_calibrate\_imu\_camera} solution (Table~\ref{tab:openitems},
item~2); completing step~3 above from a dedicated bag is the recommended
first calibration task for a new student.
\section{Publishing the report}
\label{ap:repro-report}
\begin{lstlisting}[language=bash]
cd ~/TOUT
bash tools/update_report.sh
\end{lstlisting}
\textbf{Expected:} three progress lines, then
\code{[3/3] OK, <N> pages}, then a local and a public URL.

The script runs \code{pdflatex} three times with \code{biber} in between
(needed for cross-references and the bibliography to settle), copies the PDF
into \code{web/reports/}, and \textbf{stamps the download link with a version
number}.

\noindent \textbf{Why the version stamp exists.} The Cloudflare tunnel served a
four-hour-old PDF for hours despite the origin sending
\code{Cache-Control: no-cache}, a dashboard rule overrides the origin
header, and it is not editable from this machine. A query string \emph{is} part
of the cache key, so a fresh \code{?v=} forces a miss. \textbf{Use the URL the
script prints}; a bare URL may serve a stale report, and no amount of browser
refreshing helps because the cache is at the network edge.

\noindent \textbf{Two traps inside this script, both now fixed, both worth
knowing:} \code{main.log} contains non-ASCII bytes, so \code{grep} classes it
as binary and \textbf{silently suppresses its output}, the error check saw
nothing at all until \code{-a} was added, and thirteen real errors passed while
the script reported success. And \code{pdflatex} wraps its log lines at
$\sim79$ columns, so the phrase it was grepping for the page count could be
split in two depending on document length; with \code{set -e} that killed the
script before deployment.

\section{Troubleshooting and lessons learned}
\label{ap:repro-trouble}
This section is the distilled cost of the campaign. Each entry is a real
failure, its measured signature, and its fix. They are ordered by how much
time they cost.

\subsection{The Pi overheats, and everything downstream lies}
\label{ap:repro-thermal}
\textbf{This is the most important entry in this appendix.} The Pi runs at
$86\,^\circ$C, above its soft temperature limit, so the firmware caps the CPU
at $600$\,MHz instead of $1.5$\,GHz. The consequences cascade, and none of
them look thermal:

\begin{center}
\begin{tikzpicture}[font=\scriptsize, node distance=3pt,
  s/.style={draw, rounded corners=2pt, fill=red!5, align=center, inner sep=3pt},
  a/.style={-{Latex[length=1.8mm]}, semithick}]
  \node[s] (t) {$86\,^\circ$C};
  \node[s, right=10pt of t] (f) {capped to\\$600$\,MHz};
  \node[s, right=10pt of f] (c) {CPU starved\\(load $>6$)};
  \node[s, right=10pt of c] (g) {gaps in\\\code{/imu/data\_clean}};
  \node[s, right=10pt of g] (v) {openVINS\\assertion};
  \node[s, right=10pt of v] (j) {trajectory\\jumps};
  \draw[a] (t)--(f); \draw[a] (f)--(c); \draw[a] (c)--(g);
  \draw[a] (g)--(v); \draw[a] (v)--(j);
\end{tikzpicture}
\end{center}

\noindent \textbf{How to confirm it in one command} (on the robot):
\begin{lstlisting}[language=bash]
ssh pi@$ROBOT_HOST 'vcgencmd measure_temp; vcgencmd get_throttled; \
                    vcgencmd measure_clock arm'
\end{lstlisting}
\textbf{Expected when healthy:} \code{throttled=0x0} and $\sim1.5$\,GHz. \\
\textbf{Observed at the time of writing:} \code{temp=86.2'C},
\code{throttled=0xe0006} (bits~1 and~2 set: frequency capped \emph{and}
throttled right now), \code{600} MHz.

\noindent \textbf{The measurable damage.} Count the holes in the IMU stream:
\begin{lstlisting}[language=bash]
cd ~/TOUT && source tools/robot_env.sh
source /opt/ros/noetic/setup.bash
python3 - <<'EOF'
import rospy, time
from sensor_msgs.msg import Imu
rospy.init_node('g', anonymous=True, disable_signals=True)
S=[]
rospy.Subscriber('/imu/data_clean', Imu,
                 lambda m: S.append(m.header.stamp.to_sec()), queue_size=2000)
time.sleep(20)
d=[S[i]-S[i-1] for i in range(1,len(S))]
med=sorted(d)[len(d)//2]
g=[x for x in d if x>3*med]
print('%d samples, %.1f Hz, %d gaps, %.2f s lost' %
      (len(S), 1/med, len(g), sum(g)))
EOF
\end{lstlisting}
Healthy: zero gaps. Observed: $4$ to $21$ gaps per $20$--$40$\,s, each
$\sim39$\,ms (four consecutive samples missing). Those gaps make openVINS
abort on \code{Propagator.cpp:101}, because the IMU readings between two
filter clones no longer sum to the interval between them. It respawns,
re-initialises, and the trajectory jumps.

\noindent \textbf{The fix is a fan.} \textbf{Install active cooling before
trusting any estimator measurement.} An A/B test of the same gap count before
and after cooling is the first experiment a new student should run.

\medskip
\noindent\textbf{Update, July 30: the envelope has since been measured on
both sides.} The claim in earlier editions that the software levers were
exhausted rested on an invalid test (colour was already off at driver level,
so the measurement compared a configuration with itself). The colour nodelet
has since been measured at $185\,\%$ CPU. More usefully, the failure threshold
is now bracketed rather than assumed:

\begin{center}
\small
\begin{tabular}{@{}lcc@{}}
\toprule
 & $86\,^\circ$C & $89.1\,^\circ$C \\
\midrule
Wheel odometry & $19.7$\,Hz & $0$\,Hz \\
IMU            & $82.8$\,Hz & $0$\,Hz \\
UART overruns  & $137\,000$ & $676\,751$ \\
Load average ($4$ cores) & not measured & $8.85$ \\
\code{serial\_node} & synchronised & \code{Lost sync} loop \\
\bottomrule
\end{tabular}
\end{center}

\noindent Three degrees separate a working rover from one that answers
neither manual nor autonomous commands, because \code{/serial\_node} is the
only subscriber of \code{/cmd\_vel} and it is the component that dies. If
driving stops in \emph{both} modes at once, look at the serial chain first:
that simultaneity rules out the mission logic and the web interface before you
have tested either. Recovery is a targeted cycle, not a stack restart:

\begin{lstlisting}[language=bash]
rosservice call /rosmon/start_stop "{node: 'serial_node', ns: '', action: 2}"
sleep 4
rosservice call /rosmon/start_stop "{node: 'serial_node', ns: '', action: 1}"
\end{lstlisting}

\noindent This restores the link within about twenty seconds and the launch
wrapper reapplies the \code{SCHED\_FIFO} priority automatically, but it is
a reset, not a repair. Watch the overrun counter afterwards
(\code{sudo cat /proc/tty/driver/ttyAMA}): if it keeps climbing, and at
$88\,^\circ$C it does, the collapse will recur within tens of minutes.

\subsection{The hit-and-run protocol: working usefully on a throttled Pi}
\label{ap:repro-hitrun}
The previous entry says install a fan, and that remains the right advice. This
entry is about what to do on the many days when you do not have one, because
the honest answer is not ``nothing''. A short run at $86\,^\circ$C produced
usable data; a longer one at $89\,^\circ$C produced none at all and required
manual recovery. The difference is worth managing deliberately rather than
discovering halfway through an acquisition.

\paragraph{What was actually measured, on both sides.}
With \code{enable\_color = true} and the Pi settled near $86\,^\circ$C, the
system was fully throttled (\code{throttled=0xe0006}, $600$\,MHz) and yet the
critical chain held: wheel odometry at $19.7$\,Hz, IMU at $82.8$\,Hz with
$14$ gaps per $20$\,s, serial link synchronised throughout. That is a working
rover, and it is the regime in which a short acquisition is worth attempting.

Three degrees higher the same configuration collapsed completely. Every
firmware topic dropped to $0$\,Hz, the overrun counter reached $676\,751$, and
\code{serial\_node} entered a \code{Lost sync} loop from which it did not
return unaided. \S\ref{sec:july30thermal} derives why the transition is a
cliff rather than a slope: the recovery handshake needs a clean window, and
past a threshold it stops finding one.

\begin{trapbox}{Do not read the first result as a guarantee}
``The link survived thermal throttling'' is true of the $86\,^\circ$C run and
false of the $89\,^\circ$C run. The correct statement is that the margin is
about three degrees wide with no headroom for ambient variation. Plan for the
collapse, do not assume the survival.
\end{trapbox}

\paragraph{The protocol.}
\begin{enumerate}[leftmargin=1.6em]
\item \textbf{Start cold.} Power the Pi and leave the perception stack down
  for two or three minutes. Starting at $70\,^\circ$C instead of $86\,^\circ$C
  buys most of the useful window.
\item \textbf{Open the vitals monitor before anything else}, in its own
  terminal, and leave it running for the whole session:
\begin{lstlisting}[language=bash]
python3 tools/hotrun_monitor.py        # one line every 20 s
\end{lstlisting}
  It prints temperature, throttle word, clock, IMU gap count and wheel-odometry
  rate side by side, and flags its own stop criteria.
\item \textbf{Abort on any one of three signals}, without negotiating with
  yourself: temperature above $88\,^\circ$C, more than $25$ IMU gaps per
  $20$\,s, or wheel odometry below $10$\,Hz. The third is the important one,
  because it means the serial link is going and control authority with it.
\item \textbf{Keep the run short.} Everything measured in this campaign that
  was worth keeping came from acquisitions of a few minutes. Ambition about
  duration is what turns a good run into a recovery exercise.
\item \textbf{Record to disk, not to the browser.} Use
  \code{tools/record\_trajectories.sh}, which writes a rosbag locally and is
  indifferent to the web stack. The cockpit's own buffer is cleared by any
  backend restart, and a thermal event is precisely the circumstance that
  causes one.
\end{enumerate}

\paragraph{Recovery when it does collapse.}
Cycle the serial node alone. Do not restart the stack: the July 22 lesson
(\S\ref{sec:july22network}) is that the surrounding infrastructure is almost
never the culprit, and a full restart costs minutes and destroys the run's
context.
\begin{lstlisting}[language=bash]
rosservice call /rosmon/start_stop "{node: 'serial_node', ns: '', action: 2}"
sleep 4
rosservice call /rosmon/start_stop "{node: 'serial_node', ns: '', action: 1}"
\end{lstlisting}
Expect the topics back within about twenty seconds, and expect the launch
wrapper to reapply the \code{SCHED\_FIFO} priority automatically; verify it
with \code{chrt -p <pid>} rather than assuming. Then watch the overrun counter:
\begin{lstlisting}[language=bash]
ssh pi@$ROBOT_HOST 'sudo cat /proc/tty/driver/ttyAMA' | grep -o 'oe:[0-9]*'
\end{lstlisting}
If it is still climbing, and at $88\,^\circ$C it will be, the cycle has bought
you minutes rather than fixed anything. Use them to save your data, not to
start a new run.

\paragraph{Why both modes of driving fail together, and what that tells you.}
When the serial link dies, manual and autonomous control stop at the same
instant. This looks alarming and is in fact the most useful diagnostic signal
the platform offers: the two modes share exactly one component downstream of
the mission logic, namely \code{/serial\_node}, the sole subscriber of
\code{/cmd\_vel}. A fault that removes both at once is therefore in the serial
transport, and you may discard the web interface, the command bus and the
finite-state machine without testing any of them. Confirm it in one command:
\begin{lstlisting}[language=bash]
rostopic info /cmd_vel     # publisher /leo_backend, subscriber /serial_node
rostopic hz /firmware/wheel_states   # 0 Hz => the link, not the logic
\end{lstlisting}

\subsection{A cron script that destroyed the stack every two minutes}
\label{ap:repro-watchdog}
\code{tools/leo\_watchdog.sh} runs from \code{cron} every minute and restarts
things it believes are dead. It sourced \code{/opt/ros/noetic/setup.bash} at
line~141, but its two node-presence checks call \code{rosnode} at lines~100
and~123. Under \code{cron} the \code{PATH} is minimal, so \code{rosnode} did
not exist yet: both checks failed \emph{every minute regardless of the robot},
their counters reached the 6-tick threshold, and the whole stack was torn
down, measured at six restarts in fourteen minutes.

\textbf{Why this matters to you even though it is fixed:} openVINS never had
more than two minutes to converge. Every estimator measurement taken before
the fix was made on a stack being destroyed mid-run, and had to be discarded: including a ``rotation causes divergence'' conclusion that evaporated once
the stack was stable.

\begin{itemize}[leftmargin=1.4em]
\item \textbf{Check before any measurement campaign:}
\begin{lstlisting}[language=bash]
tail -5 ~/TOUT/logs/watchdog.log
cat /tmp/leo_zombie_fails /tmp/leo_vins_fails    # both should read 0
\end{lstlisting}
\item \textbf{What else it guards, none of which was documented before
  2026-07-31.} The watchdog is not only a restart loop, and knowing what it
  does prevents you from fighting it:
  \begin{itemize}[leftmargin=1.2em]
  \item \textbf{Drowning probe on both rosbridge ports} (2026-07-13, extended
    to $9443$ on 07-31). A real WebSocket handshake against $9090$ in clear
    and $9443$ over TLS; a bridge that holds its port but has stopped
    answering is purged and relaunched. See
    \S\ref{ap:trap-zombieport}, and read the TLS warning there before
    touching it.
  \item \textbf{D455 self-healing with an escalation ladder.} A silent camera
    is first cycled through \code{rosmon}; if it stays silent, the watchdog
    issues a \emph{hardware} reset via \code{/firmware/reset\_board}. A
    persistent counter in \code{/tmp} debounces this across ticks so that one
    bad minute does not trigger a board reset. If you are power-cycling the
    camera by hand, set the maintenance flag first or the two of you will
    take turns resetting it.
  \item \textbf{Battery alert bound to the camera fault.} Past three
    consecutive silent minutes it logs the battery voltage alongside the
    fault, because a rail sagging below roughly $10.5$\,V drops the USB
    camera and presents exactly as a software failure. A camera that will not
    come back is a battery symptom until proven otherwise.
  \end{itemize}

\item \textbf{Suspend the watchdog during any deliberate manipulation:}
\begin{lstlisting}[language=bash]
touch /tmp/leo_maintenance     # watchdog does nothing while this exists
rm -f /tmp/leo_maintenance     # DO NOT FORGET THIS
\end{lstlisting}
\item \textbf{General lesson:} a script tested by hand runs with a good
  \code{PATH}. Bugs of this class only appear under \code{cron}. Test with
  \code{env -i PATH=/usr/bin:/bin bash -c '...'}.
\end{itemize}

\subsection{Wheel odometry is a reference, not the truth}
\label{ap:repro-wheels}
The mecanum wheels slip. Measured against a tape measure over $1$\,m, the
wheel odometry over-reports by $39\%$. It is still useful (it never errs by
a factor of fifty, so it catches gross estimator divergence) but it is not
ground truth and must never be quoted as accuracy.

Worse, it is \emph{asymmetric}: the correction factor is $0.9325$ turning left
and $0.8920$ turning right. Decomposing per-direction factors into a
symmetric and an antisymmetric part,
$\bar g = (g_L+g_R)/2$ and $\delta = (g_L-g_R)/2$, separates the two causes: a
\emph{sensor scale} error is sign-symmetric and lands in $\bar g$, whereas a
\emph{mechanical} defect reverses with direction and lands in $\delta$. The
wheels' $\lvert\delta\rvert$ is ten times the gyroscope's. Correcting that
with one global gain would improve one direction and worsen the other.

\subsection{When two independent sensors agree against your reference}
\label{ap:repro-indep}
A subtle trap, and a genuinely useful piece of reasoning. The first
calibration campaign found the gyroscope over-reporting rotation by $7\%$
($\bar g = 0.9329$), and the wheels over-reporting by $9\%$
($\bar g = 0.9123$). Wheel-derived heading comes from differential odometry and
\textbf{does not use the gyroscope at all}: these are physically independent
instruments. MINS and openVINS are \emph{not} independent witnesses here, since
both integrate the same gyroscope.

Two independent sensors agreeing with each other and disagreeing with the
reference points at \textbf{the reference}, not at both sensors. The
conclusion was therefore \emph{not} ``apply a $7\%$ gyro correction'' but
``re-measure with a better floor protocol''. The corresponding parameter,
\code{\textasciitilde gyro\_scale} in \code{imu\_sanitizer.py}, was added and
deliberately left at $1.0$, no correction, until three repeated passes
in each direction settle the question. It affects MINS and openVINS
simultaneously, so an error there propagates everywhere.

\subsection{Do not filter an estimator's output}
\label{ap:repro-guard}
A cautionary tale, included because it cost a full evening and the mistake is
tempting. Believing openVINS diverged during turns, the author added a node
that froze its reported position whenever the gyroscope exceeded a threshold.

It made things worse, and the measurement that proved it is simple: with the
robot verified stationary (wheel odometry $0.000$\,m over $30$\,s), raw
openVINS wandered $16.99$\,m of \emph{cumulative path} but ended
$0.021$\,m from where it started. It was oscillating around a stable point,
not drifting. The guard rejected $12\%$ of the increments of a
\emph{symmetric} oscillation, which does not remove a zero-mean signal: it
\textbf{rectifies} it. Its own counters showed $87.10$\,m of accumulated
offset for zero real motion.

\begin{itemize}[leftmargin=1.4em]
\item Distinguish \textbf{cumulative path} from \textbf{net displacement}
  before diagnosing drift. They answer different questions and only the second
  one is drift.
\item Removing samples one-sidedly biases a sum. If you must smooth, use a
  filter that \textbf{preserves the mean} (a low-pass), never a selective
  freeze.
\item Fix the estimator's inputs, not its outputs.
\end{itemize}

\subsection{Serial data loss: it is a scheduling problem}
\label{ap:repro-uart}
Symptom: thousands of checksum errors and wheel odometry collapsing from
$19$\,Hz to $1$\,Hz. It is not a cable and not the CORE2 board. The Pi's
PL011 UART FIFO overruns when \code{serial\_node} is not scheduled in time.

\begin{lstlisting}[language=bash]
ssh pi@$ROBOT_HOST 'sudo grep -o "oe:[0-9]*" /proc/tty/driver/ttyAMA'
\end{lstlisting}
\textbf{Expected:} \code{oe:0} and stable. A climbing \code{oe} counter is
overrun. The fix, already deployed, needs \emph{both} halves or the node will
not start at all: real-time priority on \code{serial\_node}
(\code{chrt -f 25}) \emph{and} a \code{LimitRTPRIO=50} systemd drop-in on
\code{leo.service} granting permission to request it. Verify with
\code{chrt -p <pid>}; expect priority~$25$.

\subsection{Configuration files that fail silently}
\label{ap:repro-yaml}
Both estimators parse their YAML with OpenCV's \code{FileStorage}
(\code{\%YAML:1.0}), which \textbf{aborts on a long comment placed after a
value on the same line}. The node then starts, subscribes to nothing, and
respawns forever: with no error naming the file. This trap was hit twice,
the second time after it had already been documented.

\begin{itemize}[leftmargin=1.4em]
\item Keep inline comments in these files \textbf{short}, or put them on their
  own line.
\item Validate before restarting anything:
\begin{lstlisting}[language=bash]
python3 -c "import cv2; f=cv2.FileStorage('estimator_config.yaml', \
cv2.FILE_STORAGE_READ); print('readable:', f.isOpened()); f.release()"
\end{lstlisting}
\item Never test a config from \code{/tmp}: relative paths inside it resolve
  differently and produce false negatives.
\end{itemize}

\subsection{Smaller traps, each of which cost real time}
\begin{itemize}[leftmargin=1.4em]
\item \textbf{\code{pgrep -f} matches your own shell.} A command containing
  the pattern appears in its own search results, so
  \code{kill \$(pgrep -f my\_node)} can kill the shell running it. Encountered
  three times. Use a bracket trick (\code{pgrep -f "my\_no[d]e"}) or filter by
  parent.
\item \textbf{\code{kill -0} lies about liveness}, because PIDs are recycled.
  Verify identity with \code{/proc/<pid>/cmdline}.
\item \textbf{Timestamps, not arrival times.} Measure rates and velocities
  from \code{header.stamp}. Arrival time mixes in network jitter and produces
  confident nonsense.
\item \textbf{Medians hide tails.} A divergence with a median of $0.69$ had a
  worst case of $19.96$ in the same data set. Report a high percentile as
  well; what ruins a trajectory lives in the tail.
\item \textbf{Path length is not comparable across sampling rates.} Per-sample
  noise inflates cumulative length as $(\sigma f/v)^2$
  (Eq.~\ref{eq:pathinflation}). Two raw lengths at $86$ and $80$\,Hz once read
  identically while the true lengths differed by $10\%$. Resample to a common
  rate first.
\item \textbf{Restart the estimators after changing calibration.} Configs are
  read once. Changing the IMU calibration mid-run once cost $34$ minutes of
  recorded data.
\item \textbf{\code{matlab -batch} needs a live licence session.} Online
  licensing fails with \code{error 5001} in batch mode when the account token
  has lapsed, even while an already-open interactive session keeps working.
  Test with \code{matlab -batch "disp('ok')"} before blaming a script.
\item \textbf{The published PDF may be cached.} The Cloudflare tunnel served a
  stale report for hours despite a correct origin header. \code{update\_report.sh}
  now stamps the download link with a version; use the link it prints, not a
  bare URL.
\end{itemize}

\subsection{The TF tree accepts only one parent per frame}
\label{ap:trap-tf}
\begin{trapbox}{Symptom: intermittent \texttt{TF\_OLD\_DATA}, poses that jump
between two plausible values, and no error anywhere}
A frame in a \code{tf2} tree has exactly \emph{one} parent. Nothing enforces
this: publish \code{base\_link} from two nodes and both transforms are
accepted, after which lookups resolve to whichever arrived last. The result is
not a crash but an estimate that silently alternates between two coherent
answers.
\end{trapbox}

\begin{figure}[htbp]
\centering
\begin{tikzpicture}[font=\scriptsize, >=Latex, node distance=9mm,
  f/.style={draw, rounded corners=2pt, inner sep=3pt, minimum width=20mm, align=center},
  bad/.style={f, draw=alertred, fill=alertred!6},
  ok/.style={f, draw=accentblue, fill=accentblue!6}]

  \node[font=\small\bfseries, align=center] at (0,1.15) {Broken: two parents};
  \node[bad] (o1) at (-1.5,0) {\code{odom}};
  \node[bad] (s1) at (-1.5,-1.5) {\code{pose\_selector}};
  \node[bad] (bf1) at (1.5,0) {\code{base\_footprint}};
  \node[bad] (bl1) at (0,-2.7) {\code{base\_link}};
  \draw[->] (o1) -- (bf1);
  \draw[->] (bf1) -- (bl1);
  \draw[->, alertred, very thick] (s1) -- (bl1);
  \node[alertred, anchor=north, align=center] at (0,-3.3)
    {two transforms, one child:\\lookups resolve to whichever arrived last};

  \node[font=\small\bfseries, align=center] at (8.2,1.15) {Fixed: strictly linear};
  \node[ok] (o2) at (8.2,0) {\code{odom}};
  \node[ok] (bf2) at (8.2,-1.35) {\code{base\_footprint}};
  \node[ok] (bl2) at (8.2,-2.7) {\code{base\_link}};
  \draw[->] (o2) -- (bf2);
  \draw[->] (bf2) -- (bl2);
  \node[accentblue, anchor=north, align=center] at (8.2,-3.3)
    {\code{base\_frame=base\_footprint}\\one parent per frame};
\end{tikzpicture}
\caption[TF tree]{The fault is structural, not numerical. On the left,
\code{base\_link} has two parents, so every lookup silently returns one of two
coherent answers. Nothing errors; the pose simply alternates. On the right,
the single-line change that restores a well-formed tree.}
\label{fig:trap-tftree}
\end{figure}

\noindent This bit twice on July 30. First, \code{pose\_selector} was
publishing with \code{base\_frame} set to \code{base\_link} while the URDF
chain already supplied \code{base\_footprint} $\to$ \code{base\_link},
giving \code{base\_link} two parents. The fix is one line in
\code{navigation\_supervision.launch}:

\begin{lstlisting}[language=XML]
<param name="base_frame" value="base_footprint"/>
\end{lstlisting}

\noindent which restores a strictly linear \code{odom} $\to$
\code{base\_footprint} $\to$ \code{base\_link} chain. Second, an early
revision of \code{carolus\_tf\_bridge.py} published the full body chain in the
beacon frame, re-parenting \code{base\_link} a second time; it now publishes
only \code{camera\_frame} $\to$ \code{carolus\_raw} plus the static anchor,
and hides the body chain behind a parameter that defaults to false.

\noindent\textbf{Check it in one command:}
\begin{lstlisting}[language=bash]
rosrun tf view_frames && evince frames.pdf
# or, faster: any frame appearing twice as a child is a fault
rosrun tf tf_monitor 2>&1 | head -30
\end{lstlisting}
Expect one parent per frame and no \code{TF\_OLD\_DATA} in the logs. Do this
before trusting any pose, because the failure is invisible in the numbers:
they stay plausible.

\subsection{The differential-drive angular multiplier, and where it is not}
\label{ap:trap-diffdrive}
The firmware applies a scalar to every angular velocity command it receives:
\begin{equation}
\omega_{\mathrm{applied}} = k_\omega \cdot \omega_{\mathrm{commanded}},
\qquad k_\omega = 1.76 .
\end{equation}
Together with the wheel radius ($0.0625$\,m) and separation ($0.358$\,m), this
is what converts \code{/cmd\_vel} into wheel setpoints, so it sits directly
in the yaw-error path: a rover that consistently over- or under-rotates by a
fixed ratio is far more likely to be showing this multiplier than a sensor
fault. Read the live values: never the file:

\begin{lstlisting}[language=bash]
rosparam get /firmware/diff_drive/angular_velocity_multiplier   # 1.76
rosparam get /firmware/diff_drive/wheel_radius                  # 0.0625
rosparam get /firmware/diff_drive/wheel_separation              # 0.358
\end{lstlisting}

\begin{trapbox}{Two ways to look for this value and find the wrong answer}
The file is \code{/etc/ros/params.yaml}: \emph{params}, plural; there is no
\code{param.yaml}. And inside it, the entire \code{diff\_drive} block
including \code{angular\_velocity\_multiplier: 1.76} is \textbf{commented
out}. Those lines document the firmware defaults; they do not set them. A
reader who greps the file finds the value inert and may conclude it is not
applied, when in fact $1.76$ is live because it is the firmware default.
Conversely, editing the commented line changes nothing until the \code{\#} is
removed \emph{and} \code{sudo systemctl restart leo} is run.
\end{trapbox}

\noindent The cockpit exposes these three values read-only in the
\emph{Firmware kinematics} panel, precisely so that they are visible during a
run without anyone being tempted to change them from a browser.

\paragraph{The order of operations for a yaw discrepancy.}
This is the part a successor most needs, because getting the order wrong
produces a rover that appears calibrated and is not. If the robot turns by a
consistently wrong amount, there are two places the discrepancy can live, and
they are not interchangeable:
\begin{equation}
\underbrace{\omega_{\mathrm{applied}} = k_\omega\,\omega_{\mathrm{cmd}}}
           _{\text{what the robot \emph{does}}}
\qquad\text{versus}\qquad
\underbrace{\hat\omega = s_g\,\omega_{\mathrm{true}}}
           _{\text{what the robot \emph{believes}}}
\label{eq:yawtwoplaces}
\end{equation}
The first is an actuation gain, set by $k_\omega =$
\code{angular\_velocity\_multiplier}. The second is a sensing gain, set by
\code{\textasciitilde gyro\_scale} in the IMU sanitiser. \textbf{Correct the
actuator first, always.} The reasoning is that $s_g$ is a property of a
physical sensor, so changing it to compensate for an actuation error stores a
lie in the one channel every downstream estimator trusts. MINS, openVINS and
the wheel odometry all consume the sanitised IMU; a wrong $s_g$ corrupts them
simultaneously and in a way no later measurement can distinguish from a real
inertial defect.

\begin{enumerate}[leftmargin=1.6em]
\item \textbf{Measure, do not infer.} Command a known rotation, for instance
  $90^\circ$, and measure the physical result with a protractor or floor
  markings. Use a deliberately unround target such as $87.5^\circ$: a round
  number invites the operator to report the number they commanded rather than
  the one they observed.
\item \textbf{Separate the two gains with one extra observation.} Compare the
  commanded angle, the physically observed angle, and the angle the gyroscope
  reported. If the robot turned too far \emph{and} the gyroscope agrees with
  the floor, the fault is in $k_\omega$. If the robot turned correctly but the
  gyroscope disagrees, the fault is in $s_g$. Only the second case justifies
  touching the sensor.
\item \textbf{Adjust $k_\omega$ multiplicatively:}
  $k_\omega^{\mathrm{new}} = k_\omega \cdot
   \theta_{\mathrm{target}} / \theta_{\mathrm{observed}}$, then re-measure.
\item \textbf{Leave \code{gyro\_scale} at $1.0$} until the ground-truth
  calibration of \S\ref{ap:repro-calibrun} has settled the outstanding
  $7$--$9\,\%$ rotation discrepancy. Until that question is answered, any
  change to $s_g$ is a guess dressed as a correction.
\end{enumerate}

\noindent A worked cautionary case is in \S\ref{sec:july29calib}: an early
calibration attempt compared the angle the robot was \emph{told} to turn
against a single entered number, and so measured how badly the platform
tracks its own command rather than how badly its sensors track reality. Those
are different quantities and only the second is a calibration.

\begin{trapbox}{A third gain, easy to forget}
A yaw error can also come from the wheel geometry itself, since
$\omega$ depends on the effective radius and the track width. Item~1 of the
open-items registry records a measured $1.65\,\%$ right/left effective-radius
bias, currently masked by a software yaw trim rather than corrected. If
$k_\omega$ needs an implausibly large change to fit, suspect the geometry
before accepting the number.
\end{trapbox}


\subsection{A TLS port that stays open and stops listening}
\label{ap:trap-zombieport}
\begin{trapbox}{Symptom: the page loads, the layout is perfect, and every
button is inert}
After a ROS master restart, the tunnel's TLS port $9443$ can remain
\emph{open}, it accepts the TCP connection and completes the WebSocket
handshake, while no longer being bound to a live rosbridge behind it.
Nothing in the browser reports an error. The interface looks connected
because, at the socket level, it is.
\end{trapbox}

\noindent\textbf{Recovery, in order of cost:} reload the page; if that fails,
connect to the local port $9090$ instead, which does not traverse the tunnel;
if that also fails, the fault is upstream of the transport and you should stop
looking at the network.

\noindent\textbf{Diagnosis before you touch anything.} Two independent checks
separate ``the page is broken'' from ``the bus is quiet'', and they are worth
running in that order:
\begin{lstlisting}[language=bash]
# 1. Does the browser's relay carry ROS traffic at all?
#    Use a topic that is ALWAYS busy. /mission/telemetry publishes ~4 Hz.
rostopic hz /mission/telemetry
# 2. Does a command sent from the page reach the bus?
rostopic echo /mission/command
\end{lstlisting}
The first check matters more than it looks. During this session the relay was
declared broken on the evidence of an idle \code{/mission/log}, which
publishes only when something happens: an empty topic proved nothing, and
the same test on \code{/mission/telemetry} showed $50$ frames in $12$\,s. A
silent topic is not a dead link. Choose a topic that is supposed to be noisy
before concluding anything.

\subsection{A checklist before any measurement campaign}
\label{ap:repro-checklist}
\begin{enumerate}[leftmargin=1.4em]
\item \code{vcgencmd measure\_temp} on the robot, below $80\,^\circ$C, and
  \code{throttled=0x0}.
\item \code{tail logs/watchdog.log}, no automatic restarts in the last hour.
\item IMU gap count over $20$\,s, zero.
\item \code{rostopic hz} on the four key topics, all at nominal rate.
\item \code{ps -o etime= -p \$(pgrep -f run\_subscribe\_msckf)}, openVINS
  older than the intended run duration, i.e.\ it has not respawned.
\item Battery on mains, or voltage noted; yaw performance depends on it.
\end{enumerate}

\noindent If any of these fails, the data you are about to record will not
survive scrutiny. The single most expensive lesson of this campaign is that a
measurement taken on an unstable stack is worse than no measurement, because
it is believed.

\bigskip\noindent\emph{The sections above are the operating manual. Those
that follow are the technical reference: where the source comes from, how it is
built, and how each configuration tree is organised.}

\section{Where to go next}
\label{ap:repro-next}
The platform is stable and instrumented; what it lacks is not more code. This
is the author's ordered recommendation, with the reasoning, so you can disagree
with it deliberately rather than by accident.

\subsection{Priority 1: re-calibrate the colour camera intrinsics}
\label{ap:next-intrinsics}
Cheapest item on this list, and it blocks the headline deliverable. The
detector runs with intrinsics calibrated at $1280\times720$ against a
$640\times480$ stream (\S\ref{sec:july30intrinsics}); the stored principal
point, $c_x = 643.47$, lies outside a $640$-pixel-wide image. A beacon at the
true image centre is reported $26.7^\circ$ off-axis. No rescaling repairs
this, because the two configurations differ in field of view and not merely in
sampling.

Run the calibration at $640\times480$, the resolution actually published,
and load the result into \code{carolus\_astrobee\_rex}. Until then, treat
every range, bearing and beacon position from this camera as a relative
indicator only. \textbf{Do not publish a dimensional metric derived from it,
including the absolute error of Priority~4.}

\subsection{Priority 2: install active cooling}
Not software. The Pi sits at $86\,^\circ$C and is frequency-capped to
$600$\,MHz whenever it is working, and that single fact taints every estimator
measurement on the platform (\S\ref{ap:repro-thermal}).

An earlier edition of this appendix stated that the software levers were
exhausted, on the grounds that disabling the colour stream gave zero thermal
improvement. That measurement was invalid: colour was already disabled at
driver level, so the test compared a configuration with itself. The colour
nodelet was subsequently measured at $185\,\%$ CPU
(\S\ref{sec:july30thermal}), making it the dominant thermal term. The lever
exists, but spending it costs the working-distance beacon detection, which
is why the July 30 session chose to keep colour and accept the risk instead.
The measured envelope is roughly three degrees wide: the serial link holds at
$86\,^\circ$C and collapses at $89\,^\circ$C. A fan is what buys back that
margin without giving up colour.

The experiment to run immediately after fitting a fan is the A/B that was never
possible: count IMU gaps over $20$\,s before and after, and count
\code{Propagator.cpp:101} aborts over a ten-minute run. If the gaps go to zero,
a whole class of ``estimator bug'' disappears, and several conclusions in
Chapter~12 should be re-measured on the healthier machine.

\subsection{Priority 3: finish the ground-truth calibration}
Three passes per direction, with the floor protocol, and settle the question
posed in \S\ref{ap:repro-calibrun}: is the $7$--$9\%$ rotation discrepancy in
the sensors or in the reference measurement? This is cheap, needs no hardware,
and is the only measurement on the platform anchored to physical reality.

Until it is settled, \code{\textasciitilde gyro\_scale} must stay at $1.0$.

\subsection{Priority 4: get an absolute reference into a driven run}
Every number in Chapter~12 except the tape-measure calibration is
\emph{mutual divergence}: it says whether MINS and openVINS agree, never which
is right. The Carolus beacon is the only drift-free anchor available, and it was
in view during neither of the two comparison runs.

Drive a loop with the beacon visible. Nothing needs to be written:
\code{plot\_trajectories.m} already computes each estimator's error against the
Carolus fix the moment a \code{\_carolus.csv} exists. That single run converts
the whole chapter from relative to absolute.

Since July 30 there is a better path than a post-hoc CSV.
\code{carolus\_tf\_bridge.py} (\S\ref{sec:july30bridge}) publishes
\code{/odom\_in\_beacon} and \code{/carolus\_in\_beacon}, both carrying the
robot's pose in the beacon frame by two independent routes, odometric and
optical. Their instantaneous difference \emph{is} the absolute error, with no
Kabsch or Umeyama fit in between, so nothing is absorbed into a fitted
transform or a scale parameter. Record both topics in the bag.

Two cautions. The bridge is deliberately in no launch file: start it by hand,
so it can never perturb a run you did not intend it to touch. And it is gated
on Priority~1: with the present intrinsics its output is self-consistent
and metrically meaningless, which is the most dangerous combination a
measurement can have, because it looks like data.

\subsection{Priority 5: the Kalibr camera--IMU calibration}
Item~2 of the open-items registry, never completed, blocked on producing
$\gtrsim1.5$\,rad/s of roll and pitch excitation with a rover that has only ever
managed $0.08$. It requires lifting and tilting the robot by hand in front of
the target. Do it after the three above, because a thermally throttled Pi will
drop frames during the very manoeuvre the calibration depends on.

\subsection{Open questions the author did not resolve}
Recorded honestly, because a successor deserves the loose ends rather than a
tidy narrative:

\begin{itemize}[leftmargin=1.4em]
\item \textbf{\code{firmware\_message\_converter} publishes an altered gyro
  bias.} It reports $6.867$ where its own source reads $3.668$. No estimator
  consumes that topic, so it was never chased, but it is a real defect and
  suggests something wrong in the conversion path.
\item \textbf{$\sim20\,000$ serial checksum errors} were logged on the Pi
  before the real-time fix. The fix cured the \emph{rate} collapse; whether it
  cured the checksum errors entirely was never re-measured.
\item \textbf{openVINS's scale disagrees with itself between runs.} One run gave
  a trajectory $10\%$ \emph{shorter} than MINS; a different run gave a
  straight-line distance $10.5\times$ \emph{longer} than reality. Both were
  measured properly. This inconsistency is unexplained and is the most
  interesting open question in the estimator work.
\item \textbf{The wheel odometry's $4.3\%$ left/right asymmetry} is
  characterised but not corrected. It is mechanical, so it needs a mechanical
  or model-level answer, not a gain.
\item \textbf{MATLAB's batch licensing} was failing at the time of writing
  (\code{error 5001}), which is why the per-panel figures in this appendix's
  companion chapter were produced by cropping rather than re-rendering. Purely
  an account/network issue, but it blocks the analysis pipeline until fixed.
\end{itemize}

\subsection{Habits worth inheriting}
\begin{enumerate}[leftmargin=1.4em]
\item \textbf{Measure before believing, and measure the thing itself.} Almost
  every wrong conclusion in this project came from measuring a proxy: arrival
  time instead of the timestamp, cumulative path instead of net displacement,
  the median instead of the tail.
\item \textbf{Write down why, not what.} Every script here has a header
  explaining the defect it works around. Six weeks later that comment is the
  only thing that distinguishes a deliberate workaround from an accident.
\item \textbf{A measurement taken on an unstable system is worse than no
  measurement}, because it will be believed. The watchdog episode
  (\S\ref{ap:repro-watchdog}) invalidated two days of estimator conclusions.
\item \textbf{Fix inputs, not outputs.} The one time this project filtered an
  estimator's output instead of its inputs, the filter became the fault
  (\S\ref{ap:repro-guard}).
\item \textbf{Keep the fallback path working.} Every convenience here has a
  manual equivalent, the cockpit calibration panel has
  \code{tools/calib\_terrain.py}, the one-click export has
  \code{bag\_to\_csv.py}. When the convenient path breaks mid-session, and it
  will, the campaign survives.
\end{enumerate}

